\section{Open parameter statements I: the published-input theorem}\label{sec:pf:open0a}

\providecommand{\openzeroaMI}{\mathrm{MI}^{*}_{2}}
\providecommand{\openzeroaHS}[1]{\|#1\|_{\mathrm{HS}}}
\providecommand{\openzeroahs}{\operatorname{dev}}
\providecommand{\openzeroamfL}{\operatorname{mfL}}
\providecommand{\openzeroakfk}{\operatorname{kfk}}
\providecommand{\openzeroaeD}{e_{\mathrm{D}}}
\providecommand{\openzeroacj}{\operatorname{cj}}
\providecommand{\openzeroaLk}{\mathrm{Lk}}
\providecommand{\openzeroaEv}{\mathrm{Ev}}
\providecommand{\openzeroaTrue}{E^{\mathrm{true}}}
\providecommand{\openzeroaram}{M^{\mathrm{ram}}}
\providecommand{\openzeroaRow}{\mathcal{R}_0}
\providecommand{\openzeroaNV}{\mathrm{NV}}
\providecommand{\openzeroaSC}{\mathrm{SC}}

This section proves the theorem of the formal development \cite{LeanSource2026} which carries the multiplicity
statement among the twelve open parameter statements of \cref{thm:families:open}: under thirteen classical
results on unitary groups and a diagonal row hypothesis, the \emph{multiplicity identity} $\openzeroaMI$ holds,
eventually in the parameters, at every bundle linked to a good point.

\subsection{Introduction}\label{sec:pf:open0a:intro}

\paragraph{The result.}
A bundle $b$ over a member of a family carries an \emph{indefinite realisation}, and the realisation determines a
\emph{coefficient complex} $X_b$: Hilbert spaces of theta images of rows, a Hilbert space of automorphic forms
containing period vectors of columns, and algebraic \emph{column pairings} between columns and rows with values in
the coefficient field. The identity $\openzeroaMI(X_b)$ (\cref{pf:open0a:def:mistar}) asserts that, after
normalisation by explicit unit-like constants, the column pairings are inner products
\[
\zeta(u,t)\,\iota_0\bigl(\mathrm{Pt}(u,t)(\varphi,F)\bigr)
=a(u)\,\bigl\langle s_{EE}\bigl(V\,\Theta_t(\cdots F)\bigr),\ P_u\bigl(U(u)\,\iota_0(\varphi)\bigr)\bigr\rangle
\]
between the theta image of the row $F$ and the period vector of the column $\varphi$. The main result of the
section, \cref{pf:open0a:main}, asserts this identity, in the eventual sense of the row language of
\cref{def:families:rows}, at all bundles linked to good points, for every parameter $p$. It is applied in the proof
of \cref{thm:families:open}, where the diagonal row hypothesis is supplied by the convergent Siegel--Weil formula and
the remaining hypothesis of \cref{pf:open0a:main} by the construction of the points.

\paragraph{The argument.}
The proof has three parts.
\begin{enumerate}[label=(\roman*)]
\item \emph{True packages} (\cref{sec:pf:open0a:points}). Eventually, every bundle linked to a good point satisfies
ten \emph{point conditions}, and together with its bundle witness these make the real package of the bundle
\emph{true}: a list of family, level and character conditions holds.
\item \emph{Non-vanishing of periods} (\cref{sec:pf:open0a:punz}). For every word $u$ some ramified vector has a
non-zero period. This is a dimension count: the kernel of the period map $P_u$ has dimension at most the
Hilbert--Schmidt defect $\openzeroaHS{K_u-1}^2$ of the normalised Gram operator $K_u$ of $P_u$; the diagonal row
hypothesis bounds this defect by a small multiple of $\dim C_u$, while the ramified vectors occupy a proportion
$(1-1/(T_u-1))^f\ge1-f/(T_u-1)$ of $C_u$.
\item \emph{The multiplicity identity for true packages} (\cref{sec:pf:open0a:mi}). For a true package with
non-vanishing periods, $\openzeroaMI$ is derived from: surjectivity of the theta maps along chains (Harish-Chandra
theory, Howe's unramified computation and the doubling method); family conditions along iterated seesaws (strong
approximation and the theory of discrete series); seesaw transport relations, which transfer linear relations
between top pairings to linear relations between bottom periods; a central element of the finite adelic group which
pins the bottom data; and an \emph{anchoring} argument, in which the non-vanishing of periods provides one non-zero
pairing per word and linear algebra propagates the relations to all pairings.
\end{enumerate}

\paragraph{What is proved in full.}
The assembly of the main theorem, the non-vanishing of periods with all its numerical lemmas, the served condition,
the joint condition from the seesaw transport relations, the central shift, the anchoring and normalisation
arguments and the final construction of the data of $\openzeroaMI$ are proved in full. The ingredients with long
formal proofs (the size estimates at linked bundles, the leading reads, the surjectivity along chains, the geometric
families, the column choices along seesaws, the central pin, the seesaw transport relations, the translated
evaluations, the pairing form and the bottom relation) are proved in \cref{sec:pf:open0a:deferred} at the level of
their main steps; the deeper ingredients are stated precisely and cited.

\paragraph{Organisation.}
\Cref{sec:pf:open0a:setup} recalls the inputs and the row language and states the main theorem.
\Cref{sec:pf:open0a:reduction} reduces it to four results, which are proved in
\cref{sec:pf:open0a:punz,sec:pf:open0a:points,sec:pf:open0a:mi}, using the results of
\cref{sec:pf:open0a:deferred}. \Cref{sec:pf:open0a:remarks} collects remarks on the formal statements.

%==============================================================================
\subsection{The setting and the main theorem}\label{sec:pf:open0a:setup}

\subsubsection{The thirteen published inputs}
The theorem assumes thirteen propositions of the development, each the formal statement of a classical theorem;
they are fields of the package $P$ of \cref{sec:families:theta} and are proved in the development
(\cref{sec:literature}). We list them with the form in which they are stated and refer to them as (P1)--(P13).
\begin{enumerate}[label=(P\arabic*),leftmargin=*]
\item\label{pf:open0a:P1} \emph{Finiteness} ($\sfname{HarishChandraFinite}$): for a datum $D$, a hermitian space
$V$, an open subgroup $U$ of the finite adelic points of $\Uni(V)$, archimedean frames, finite-dimensional
continuous representations $\rho_\tau$ of the maximal compact subgroups and a system of eigenvalues $c$, the span of
the automorphic forms on $\Uni(V)$ of level $U$, $K_\infty$-type $\rho$ and infinitesimal character $c$ is
finite-dimensional \cite{HarishChandra1968,Borel1969}.
\item\label{pf:open0a:P2} \emph{Rapid decay} ($\sfname{CuspRapidDecay}$): smooth cuspidal automorphic functions of
uniform moderate growth on $\Uni(V)(F)\backslash\Uni(V)(\A_F)$ are of rapid decay
\cite{HarishChandra1968,MoeglinWaldspurger1995}.
\item\label{pf:open0a:P3} \emph{Discrete series} ($\sfname{DiscreteSeriesCover}$): for $\Uni(p,q)$, $p+q\ge1$, and
its finite covers there is a family $\pi_\lambda$ of discrete series representations indexed by Harish-Chandra
parameters $\lambda$, with the expected $K$-types, such that every discrete series representation is unitarily
equivalent to some $\pi_\lambda$ \cite{HarishChandra1966}.
\item\label{pf:open0a:P4} \emph{Harish-Chandra modules}: Harish-Chandra modules of irreducible unitary
representations of $\Uni(p,q)$ and their covers, and analyticity of $K$-finite vectors \cite{HarishChandra1953}.
\item\label{pf:open0a:P5} \emph{Casselman--Wallach globalisation}: a $(\mathfrak g,K)$-map from the $K$-finite
vectors of an irreducible unitary representation of $\Uni(p,q)$ into a smooth Fr\'echet representation of moderate
growth extends uniquely to a continuous intertwiner on smooth vectors, with closed image and control of Sobolev norms
\cite{Casselman1989,Wallach1992}.
\item\label{pf:open0a:P6} \emph{Howe's unramified computation} ($\sfname{HoweUnram}$): for an unramified local dual
pair of unitary groups and a representation of the pair with a vacuum vector $s$ generating the Heisenberg action,
the vectors fixed by the hyperspecial subgroup of one member lie in the span of the translates of $s$ by the other
member \cite{Howe1979}.
\item\label{pf:open0a:P7} \emph{Weil representations}: global Weil representations of unitary dual pairs with
prescribed splitting characters exist \cite{Kudla1994}.
\item\label{pf:open0a:P8} \emph{Weil's Siegel formula}: convergence of theta integrals and the identity of
Siegel--Weil type in Weil's range of absolute convergence \cite{Weil1965}.
\item\label{pf:open0a:P9} \emph{Siegel--Weil formula}: the Siegel--Weil formula for the doubled space in the
convergent range \cite{Weil1965,KudlaRallis1988,Ichino2004}.
\item\label{pf:open0a:P10} \emph{Doubling method}: continuation of Siegel Eisenstein series on $\Uni(V\oplus-V)$
with singularities of finite order, the basic identity of the doubling zeta integral and its continuation
\cite{GelbartPiatetskiShapiroRallis1987,PiatetskiShapiroRallis1987}.
\item\label{pf:open0a:P11} \emph{Unramified doubling integral}: for $\GL_n$ it equals
$L(s+\frac12,\alpha)L(s+\frac12,\alpha^{-1})$ times the normalising factor \cite{GelbartPiatetskiShapiroRallis1987}.
\item\label{pf:open0a:P12} \emph{Classification of discrete series}: for each $n$ there are unitary
square-integrable irreducible models $H(\varepsilon,\lambda)$ for regular integral $\lambda$, with infinitesimal
character $\lambda$, a unique lowest $K$-type $\lambda+\rho_n-\rho_c$ and all $K$-types of the form
$\lambda+\rho_n-\rho_c+(\text{sum of positive roots})$, and every irreducible square-integrable unitary
representation is unitarily equivalent to some $H(\varepsilon\circ s,\mu)$ \cite{HarishChandra1966,Knapp1986}.
\item\label{pf:open0a:P13} \emph{Strong approximation} for $\mathrm{SU}(J)$, $J$ a non-degenerate hermitian matrix
of size $\ge2$ over a CM field, with respect to a finite set $S$ of places such that $\mathrm{SU}(J)(F_S)$ is not
compact \cite{Kneser1966,PlatonovRapinchuk1994}.
\end{enumerate}

\subsubsection{Points, bundles and eventual statements}
We use the row language of \cref{def:families:closed,def:families:rows}, which we make precise as far as needed.
A \emph{context} $K$ fixes a scale $K.s=(x,\varepsilon,\dots)$, a system $K.C$ of family constants and further reals
$c_D,c_h$. A \emph{point} $\beta$ in the context $K$ is an output of the member construction at scale $K.s$; it
determines a datum $D_\beta$ and an object $O_\beta$ over $D_\beta$ (size data, theta and twist data). A \emph{slot}
is a predicate on points. For a datum $D$, an object $O$ over $D$ and $c\in\R$, a \emph{bundle} $b$ over $(D,O,c)$
consists of a reciprocity record $b.R$, real spectral ties $b.S$ (which contain a Hida setting, i.e.\ an automorphic
instance of a unitary group with level and archimedean data), a row-side family $b.P$, square-integrable columns
$b.C$ and a real theta package $b.B$, which contains an indefinite realisation. A \emph{linkage} is a relation
$\openzeroaLk(K,\beta,c,b)$ between points and bundles over them.

\begin{definition}[Eventual statements]\label{pf:open0a:def:rowE}
A \emph{parameter tuple} $\kappa=(\varepsilon,C,c_D,c_h,C_5,c_6,C_b)$ consists of a system $C$ of family constants
and six reals; it is \emph{positive} if $C_5,c_6,C_b>0$. A context $K$ is \emph{of type} $\kappa$ if $K.C=C$,
$K.c_D=c_D$, $K.c_h=c_h$, $K.s.\varepsilon=\varepsilon$ and the goodness condition of the scale with respect to
$(C,C_5,c_6,C_b,c_D,c_h)$ holds.
\begin{enumerate}[label=(\alph*)]
\item A property $q$ of parameter tuples holds \emph{for large parameters} if there is $\varepsilon_0>0$ such that
for all reals $\varepsilon,C_8,C_5,c_6,C_b$ with $0<\varepsilon\le\varepsilon_0$, $\varepsilon<1/52$ and $C_8\ge1$
there is a staged threshold $K_0$ such that for every standard system $C$ of family constants with prescribed constant
$C_{\mathrm{pde}}=C_8$ which is large with respect to $K_0$ there is $c_1$ such that for every $c_D\ge c_1$ there is $c_2$
such that $q(\varepsilon,C,c_D,c_h,C_5,c_6,C_b)$ holds for every $c_h\ge c_2$.
\item A property $p(\kappa,K)$ holds \emph{eventually in contexts} if the following property of $\kappa$ holds for
large parameters: if $\kappa$ is positive, there is $x_0$ such that $p(\kappa,K)$ holds for every context $K$ of
type $\kappa$ with $K.s.x\ge x_0$.
\item For a slot $Q$, a linkage $\openzeroaLk$ and a property $E(\kappa;D,O,c,b)$ of parameter tuples and bundles,
$\openzeroaEv_{Q,\openzeroaLk}(E)$ is the statement that the following holds eventually in contexts: for every
point $\beta\in Q$ in context $K$, every $c\in\R$ and every bundle $b$ over $(D_\beta,O_\beta,c)$ with
$\openzeroaLk(K,\beta,c,b)$, the property $E(\kappa;D_\beta,O_\beta,c,b)$ holds. We then say that $E$ holds
\emph{eventually at the bundles of} $(Q,\openzeroaLk)$.
\end{enumerate}
\formal{GlobalTheory.Fc.ParEv}\ \formal{GlobalTheory.Fc.EvP}\ \formal{GlobalTheory.Fc.RowE}
\end{definition}

\begin{lemma}[Calculus of eventual statements]\label{pf:open0a:rowEcalc}
Let $Q$ be a slot and $\openzeroaLk$ a linkage.
\begin{enumerate}[label=(\alph*)]
\item If $\openzeroaEv_{Q,\openzeroaLk}(E)$ and $\openzeroaEv_{Q,\openzeroaLk}(E')$, then
$\openzeroaEv_{Q,\openzeroaLk}(E\wedge E')$.
\item If $\openzeroaEv_{Q,\openzeroaLk}(E)$, and $E(\kappa;D_\beta,O_\beta,c,b)$ implies
$E'(\kappa;D_\beta,O_\beta,c,b)$ whenever $K$ is of type $\kappa$, $\beta\in Q$ is a point in context $K$ and
$\openzeroaLk(K,\beta,c,b)$, then $\openzeroaEv_{Q,\openzeroaLk}(E')$.
\end{enumerate}
\formallabel{GlobalTheory.Fc.RowE.and}\ \formallabel{GlobalTheory.Fc.RowE.imp}
\end{lemma}

\begin{proof}
(b) is immediate from the definitions: the same thresholds work. For (a) it suffices to show that two properties
which hold eventually in contexts hold simultaneously eventually in contexts. Given $\varepsilon_0,\varepsilon_0'$,
take $\min(\varepsilon_0,\varepsilon_0')$; given the staged thresholds $K_0,K_0'$, take their componentwise maximum
(a system of constants which is large with respect to it is large with respect to both); and take the maxima of the
two lower bounds $c_1$, of the two lower bounds $c_2$ and of the two lower bounds $x_0$.
\end{proof}

\subsubsection{Parameters and linkage}
Let $h_L$ be a proof of $\sfname{RecTExists}$ (\cref{def:rectexists}). A \emph{parameter} $p$ for $h_L$ consists of a
position, a norm-choice extension, a spectral extension and a geometric predicate (\cref{def:families:rows}); it
determines a system of slots and the slot $\mathcal G_p$ of \emph{good} points ($\sfname{IsGoodE}(p)$ in
\cref{def:families:rows}). A \emph{weighted point} $q$ of $p$ in the context $K$ consists of a base, a frame, twisting
data, column data $w$ at a later stage $Z$, an object $O_q$ over the datum $D_q=w.D(Z)$, and a \emph{witness}, a long
record of compatibilities which contains the \emph{reading data} of the point. A \emph{bundle witness} for a bundle $b$
over $(D_q,O_q,c)$ is a record relating $b$ to $q$: a twisted bundle input for $b$, the compatibility of the
reciprocity record of $b$ with $q$, the identification of the spectral ties of $b$ with the saturation of $q$, the
identification of column weights and of division data, the selected side of the real record, and the compatibility
of the column characters.

\begin{definition}[The linkage of a parameter]\label{pf:open0a:def:link}
For a parameter $p$, a point $\beta$ in context $K$, $c\in\R$ and a bundle $b$ over $(D_\beta,O_\beta,c)$, the
relation $\openzeroaLk_p(K,\beta,c,b)$ holds if there are a weighted point $q$ of $p$ in context $K$ and a bundle $b'$
over $(D_q,O_q,c)$ admitting a bundle witness, such that $(D_\beta,O_\beta,b)=(D_q,O_q,b')$ as dependent triples.
\formal{GlobalTheory.Mem.Top.Bl.BundleLinkage}
\end{definition}

\subsubsection{The coefficient complex and the multiplicity identity}
Let $\mathcal{K}$ be a number field with a complex conjugation, $v_0$ an infinite place of $\mathcal{K}$ and
$\iota_0\colon\mathcal{K}\to\C$ the corresponding embedding. For a finite place $w$ of $\mathcal{K}$ and
$k\in\mathcal{K}^\times$ let $v_w(k)\in\Z$ be the exponent of $w$ in the fractional ideal $k\OO_{\mathcal{K}}$, and
$v_w(0)=0$. For $Y\in\GL_a(\C)$ and $y\in\C^a$ put $\|y\|_Y=\|Y^{-1}y\|$ (Euclidean norm). Inner products of Hilbert
spaces are linear in the second variable.

\begin{definition}[Coefficient complexes]\label{pf:open0a:def:cx}
A \emph{coefficient complex} $X$ over a base with a set $\mathcal T$ of \emph{tuples}, a base tuple
$\mu_0\in\mathcal T$, and an ambient assigning $\mathcal{K}$-vector spaces $F_R(t)$ to $t\in\mathcal T$, consists of
(among further data not used here): finite-dimensional Hilbert spaces $F_C(t)$ with $\iota_0$-semilinear maps
$j_R(t)\colon F_R(t)\to F_C(t)$; a finite-dimensional Hilbert space $E$ with linear maps $\Theta_t\colon F_C(t)\to E$
and a subspace $\openzeroaTrue\subseteq E$, the \emph{true part}; and a map $\mu$ from the finite places of
$\mathcal{K}$ to $\mathcal T$. The \emph{face} of $X$ (attached to the sub-datum $t\mapsto
j_R(t)^{-1}(\operatorname{im}\Theta_t^*)$) provides: integers $n_V(t)$ (rows at $t$ are vectors
$F\in\mathcal{K}^{n_V(t)}$); a set $\mathrm{Mem}$ of \emph{members} $u$ with integers $a(u)$ (columns of $u$ are
vectors $\varphi\in\mathcal{K}^{a(u)}$), \emph{column metrics} $X_c(u,v)\in M_{a(u)}(\C)$ at the infinite places $v$,
and $\mathcal{K}$-bilinear \emph{column pairings}
$\mathrm{Pt}(u,t)\colon\mathcal{K}^{a(u)}\times\mathcal{K}^{n_V(t)}\to\mathcal{K}$; and a Hilbert space $H$ with
linear \emph{period maps} $P_u\colon\C^{a(u)}\to H$.
\formal{Formalism.IndefR.Cx.CxData}
\end{definition}

\begin{definition}[The multiplicity identity $\openzeroaMI$]\label{pf:open0a:def:mistar}
Let $X$ be a coefficient complex, $e_V(t)\colon\mathcal{K}^{n_V(t)}\to F_R(t)$ $\mathcal{K}$-linear maps (the
\emph{row embeddings}) and $s_{EE}\colon E\to H$ a conjugate-linear map. Then $\openzeroaMI(X,e_V,s_{EE})$ holds if
there are a linear map $V\colon E\to E$, constants $a(u)\in\C$, matrices $U(u)\in M_{a(u)}(\C)$ and constants
$\zeta(u,t)\in\C$ ($u\in\mathrm{Mem}$, $t\in\mathcal T$) such that:
\begin{enumerate}[label=(M\arabic*)]
\item\label{pf:open0a:M1} $V(\openzeroaTrue)\subseteq\openzeroaTrue$ and $\|Ve\|=\|e\|$ for $e\in\openzeroaTrue$;
\item\label{pf:open0a:M2} $\|U(u)y\|_{X_c(u,v_0)}=\|y\|_{X_c(u,v_0)}$ for all $u$ and $y\in\C^{a(u)}$;
\item\label{pf:open0a:M3} $\zeta(u,t)\ne0$ for all $u,t$; $\zeta(u,\mu_0)=1$; and for every $u$, every finite place $w$
and every $k\in\mathcal{K}$ with $\iota_0(k)=\zeta(u,\mu(w))$ one has $v_w(k)=0$;
\item\label{pf:open0a:M4} $a(u)\ne0$ for all $u$;
\item\label{pf:open0a:M5} for all $u,t$, $\varphi\in\mathcal{K}^{a(u)}$ and $F\in\mathcal{K}^{n_V(t)}$,
\begin{equation}\label{pf:open0a:eq:MI}
\zeta(u,t)\,\iota_0\bigl(\mathrm{Pt}(u,t)(\varphi,F)\bigr)
=a(u)\,\bigl\langle s_{EE}\bigl(V\,\Theta_t(j_R(t)(e_V(t)F))\bigr),\ P_u\bigl(U(u)\,\iota_0(\varphi)\bigr)\bigr\rangle ,
\end{equation}
where $\iota_0(\varphi)\in\C^{a(u)}$ is taken componentwise.
\end{enumerate}
\formal{Formalism.IndefR.Cx.CxData.MIstar2}
\end{definition}

For a bundle $b$ we write $X_b$ for the coefficient complex of its indefinite realisation, $e_{V,b}$ for the row
embeddings of the face of the realisation and $s_{EE,b}$ for its conjugate-linear map; the number field is the
coefficient field $\mathit{LL}$ of the datum, with the embedding $\iota_0$ of the realisation.

\subsubsection{Words, periods and the diagonal row hypothesis}\label{sec:pf:open0a:words}
Let $D$ be a datum with ranks $k,m,f,r,r_S,\dots$; recall that $m\ge2$, $k\ge m+3\ge5$ and $f=2n$ with $n\ge1$, so
that $f\ge2$. The spectral data of $D$ attach to every word $u$ a finite-dimensional Hilbert space $C_u$ of
automorphic forms, a Hilbert space $H$ (independent of $u$), a linear \emph{period map} $P_u\colon C_u\to H$ and a
real constant $c_0(u)$ (a ratio of volumes and dimensions). The member of $u$ provides a prime $T_u$, the \emph{level
prime}, the space $M_0(u)$ of forms of level $T_u$ with a linear isomorphism $\iota_u\colon C_u\to M_0(u)$, and the
\emph{ramified subspace} $\openzeroaram(u)\subseteq M_0(u)$, the common kernel of the norm maps at the places above
$T_u$. Put
\[
d_0(u)=\dim C_u=\dim M_0(u),\qquad d_U(u)=\dim\openzeroaram(u),\qquad K_u=(c_0(u)^2)^{-1/2}\,P_u^*P_u,
\]
and let $\openzeroahs(u)=\openzeroaHS{K_u-\mathrm{id}}^2$ be the \emph{diagonal defect} (Hilbert--Schmidt norm). The
member of $u$ also provides reals $\Gamma_u>1$, $D_*(u)>0$ and an integer $d_u^*\ge1$; with
\begin{equation}\label{pf:open0a:eq:Tstar}
X_u=\bigl(D_*(u)/d_u^*\bigr)^{1/(fk)},\qquad T_u^*=X_u\,\Gamma_u^{-2/(fk)}
\end{equation}
the level prime lies in the window $T_u^*\le T_u\le(1+\frac1{fk})\,T_u^*$; moreover $D_*(u)=D_*$ for all $u$, where
$D_*>0$ is the main term of $D$ (\cref{thm:theorem-S}).

\begin{definition}\label{pf:open0a:def:diag}
Let $D$ be a datum and $u$ a word. Put $\openzeroamfL(L)=m^5fL$, $\openzeroakfk=kf\log k$ and
$\openzeroaeD(L,C)=8e^{-C\openzeroamfL(L)}+6e^{-\openzeroakfk/2}$ for $L,C\in\R$.
\begin{itemize}
\item The \emph{rank identity} at $u$: $d_0(u)=d_u^*\,(T_u^k-1)^f$.
\item The \emph{two-term bound} $\mathrm{TT}_u(A,x)$: $\openzeroahs(u)\le A\,d_0(u)^2/D_*+x\,d_0(u)$.
\item The \emph{diagonal bound} $\Delta_u(L,C)$: $\mathrm{TT}_u\bigl(4(e^{C\openzeroamfL(L)}+1),\
4e^{-\openzeroakfk/2}+2e^{-\openzeroakfk}\bigr)$.
\item \emph{Smallness} $\mathrm{Sm}_u(L,C)$: $d_0(u)\le D_*\,e^{-2C\openzeroamfL(L)}$.
\item The \emph{defect bound} $\mathrm{Df}_u(e)$: $\openzeroahs(u)\le e\,d_0(u)$.
\item The \emph{ramified share} $\mathrm{Rs}_u$: $d_U(u)=d_0(u)\,(1-\frac1{T_u-1})^f$; the \emph{ramified bound}
$\mathrm{Rb}_u(\delta)$: $d_U(u)>0$ and $(1-\delta)\,d_0(u)\le d_U(u)$.
\item \emph{Non-vanishing of periods} $\openzeroaNV(D)$: for every word $u$ there is $\varphi\in\openzeroaram(u)$ with
$P_u(\iota_u^{-1}\varphi)\ne0$.
\end{itemize}
\formal{GlobalTheory.Core.Rk.SDiag}\ \formal{GlobalTheory.Core.Rk.DiagS}\ \formal{GlobalTheory.MiTop.PuNZ}
\end{definition}

The diagonal row hypothesis is formulated with the \emph{standing size conditions}
$\openzeroaSC(D;L,\mathbf C)$, $\mathbf C=(C',C_l,C_g,c_b,C_u,c_T,C_T)$: the conjunction of the ties of $D$, the
mid-volume formula for the dimensions of the words, a list of explicit inequalities between the ranks, $L$, the root
discriminant, the type primes $g_1,g_2$, the Harish-Chandra parameters, the level primes and the constants
$\mathbf C$ (the \emph{size inequalities}), a list of interlacing and congruence conditions, and the trace frame. Of
these we use only
\begin{equation}\label{pf:open0a:eq:HDused}
L>1,\quad m^4\le f\le k,\quad (ke^L)^8\le g_1,\quad (ke^L)^8\le g_2,\qquad
T_u\equiv1\ \bmod\ (g_1g_2)^2,\quad d_0(u)>0,
\end{equation}
the first four from the size inequalities, the last two from the congruence conditions. A datum is \emph{true} if it
is the datum $I.\mathrm{rec}(l,Z)$ produced by the member construction from a base, a pre-datum, a construction $I$,
laws $l$ and later data $Z$, and a \emph{world} over $(I,l,Z)$ exists; it is \emph{true in the strong sense} if
moreover the Hida setting of some such world satisfies the identity condition of its coordinates.

\begin{definition}[The diagonal row hypothesis]\label{pf:open0a:def:rowdiag}
For a predicate $E(D,L,C')$, the \emph{diagonal row hypothesis} $(\mathrm{DR}_E)$ is the statement: for all reals
$C_g,c_b,C_u,c_T,C_T$ with $c_b>0$ and $c_T>0$ there are reals $C\ge1$, $C'_0$, $C_{l,0}$ such that for all
$C'\ge C'_0$, $C_l\ge C_{l,0}$, every datum $D$, every $L\in\R$ and every word $u$ of $D$: if $D$ is true in the strong
sense, $E(D,L,C')$ and $\openzeroaSC(D;L,C',C_l,C_g,c_b,C_u,c_T,C_T)$, then $\Delta_u(L,C)$. We use it with
$E=E_{\mathrm{diag}}$, where $E_{\mathrm{diag}}(D,L,C')$ means: the rank identity holds at every word, and $D$ has size
data $L_2$ with scale parameter $L_2.L=L$ and constant $C_\Gamma(L_2)$ satisfying $1\le C_\Gamma(L_2)$ and
$8C_\Gamma(L_2)\le C'$. Every size data $L_2$ of $D$ satisfies
\begin{equation}\label{pf:open0a:eq:Gam}
\log\Gamma_u=C_\Gamma(L_2)\,m^5f\,L_2.L\qquad\text{for all words }u.
\end{equation}
\formal{GlobalTheory.Core.Rk.RowDiag}
\end{definition}

\subsubsection{The main theorem}

\begin{theorem}[The published-input theorem]\label{pf:open0a:main}
Assume \ref{pf:open0a:P1}--\ref{pf:open0a:P13} and $(\mathrm{DR}_{E_{\mathrm{diag}}})$. Let $h_L$ be a proof of
$\sfname{RecTExists}$ and $p$ a parameter for $h_L$, and assume that for every context $K$ and every weighted point $q$
of $p$ in context $K$ the reading data of $q$ admit a reading structure. Then $\openzeroaMI(X_b,e_{V,b},s_{EE,b})$ holds
eventually at the bundles of $(\mathcal G_p,\openzeroaLk_p)$: eventually in contexts, for every good point $\beta$,
every $c\in\R$ and every bundle $b$ over $(D_\beta,O_\beta,c)$ linked to $\beta$, the coefficient complex of $b$
satisfies the multiplicity identity.
\formallabel{GlobalTheory.T.main_theorem_published_inputs_variant_alt_alt}\
\formallabel{GlobalTheory.miStarRowClosing_alt_alt}
\end{theorem}

Here a \emph{reading structure} on the reading data of a weighted point is a record relating the reading maps of the
point to the ambient spaces of its column data (agreement of the base, moved and twisted readings and of the
coordinates), together with vanishing conditions at unused tuples, a condition on the length of a Tate--Shafarevich
group, the laws, canonicity, positivity of the ranks and standardness of the core of the point.


%==============================================================================
\subsection{Reduction to four results}\label{sec:pf:open0a:reduction}

\subsubsection{Real packages}
Let $D$ be a datum and $\mathrm{ex}$ an extra datum over $D$ (levels, theta and twist data; it is a component of every
object $O$ over $D$). A \emph{real package} $W$ over $\mathrm{ex}$ consists of a reciprocity record $W.R$, real spectral
ties $W.S$ over $W.R$ (with a Hida setting), a row-side family $W.P$, column data and period data for every sub-datum
of the ambient of $W.P$, scale data, a number field with an embedding into $\mathit{LL}$, and an indefinite
realisation $W.Q$. Its coefficient complex $X_W$, row embeddings $e_{V,W}$ and conjugate-linear map $s_{EE,W}$ are
those of the realisation $W.Q$. A bundle $b$ over $(D,O,c)$ determines the real package $\mathrm{pk}(b)$ over
$O.\mathrm{ex}$, with components $b.R$, $b.S$, $b.P$, the column faces and periods of $b.C$, the scale data determined
by $b.P$, the division data of $b$ and $c$, and the number field and realisation of $b.B$. By construction
$X_{\mathrm{pk}(b)}=X_b$, $e_{V,\mathrm{pk}(b)}=e_{V,b}$ and $s_{EE,\mathrm{pk}(b)}=s_{EE,b}$.

\begin{definition}[Point conditions; true packages]\label{pf:open0a:def:true}
\begin{enumerate}[label=(\alph*)]
\item A bundle $b$ over $(D,O,c)$ satisfies the \emph{point conditions} $\mathcal P(b)$ if: a join of the
realisation of $b$ with $b.R$ and $b.S$ exists; the decomposition point condition of $b.P$ holds; frames of the
realisation exist; the realisation extends; the existence condition of its real datum holds; its primitivity rows
hold and the types of $b.P$ are separated; the finiteness condition of the spectral data of $b.S$ holds; the
Harish-Chandra parameters of the Hida setting of $b.S$ have gaps $\ge100\,r^4$ between entries of different signs
(the \emph{gap condition}); the node gap condition of $D$ holds (\cref{sec:pf:open0a:points}); and $r_S\ge1$.
\item A real package $W$ satisfies the \emph{package conditions} $\mathcal A(W)$ if: it satisfies the \emph{package
hypotheses} $\mathrm{PH}(W)$, which contain the \emph{bottom hypotheses} $\mathrm{BH}(W)$ and the \emph{level
hypothesis} $\mathrm{LH}(W)$ (\cref{sec:pf:open0a:mi}); it is \emph{valid} (the construction is admissible, pins
exist, it is twisted, standard columns exist, the characters $\mu_0$ match those of the construction, and a
\emph{leading read} exists); it satisfies the \emph{geometric conditions} (the characters $\xi_t$ and the prime sets
$Q_t$ match those of the construction, and leading reads exist at the identity); the standard twist condition and the
level condition hold; and a recursion of the realisation exists.
\item $W$ is \emph{true} if $D=I.\mathrm{rec}(l,Z)$ for a member construction (and $\mathrm{ex}$ lies over this
datum) and $\mathcal A(W)$ holds.
\end{enumerate}
\formal{GlobalTheory.Fam.BundlePointConditions}\ \formal{GlobalTheory.Fam.AllFamilyConditions}\
\formal{GlobalTheory.Fam.TruePackAlt}
\end{definition}

\subsubsection{The four results}
The main theorem follows from four results, proved in \cref{sec:pf:open0a:points,sec:pf:open0a:punz,sec:pf:open0a:mi}.

\begin{proposition}[Point conditions at linked bundles]\label{pf:open0a:points}
Assume \ref{pf:open0a:P1}--\ref{pf:open0a:P5}. For every $h_L$ and every parameter $p$, the point conditions
$\mathcal P(b)$ hold eventually at the bundles of $(\mathcal G_p,\openzeroaLk_p)$.
\formallabel{GlobalTheory.T.point_lines_of_cusp_rapid_decay}\
\formallabel{GlobalTheory.Fam.PtRowsAt.good_linked_ptR_of_cuspRapidDecay}
\end{proposition}

\begin{proposition}[True packages from bundle witnesses]\label{pf:open0a:truepack}
Let $q$ be a weighted point of $p$ in context $K$ and $b$ a bundle over $(D_q,O_q,c)$ with a bundle witness. If the
reading data of $q$ admit a reading structure and $\mathcal P(b)$ holds, then the real package $\mathrm{pk}(b)$ is
true.
\formallabel{GlobalTheory.T.true_pack_of_bundle_witness}\ \formallabel{GlobalTheory.Fam.truePack_of_bWit_alt}
\end{proposition}

\begin{proposition}[Non-vanishing of periods]\label{pf:open0a:punz}
Assume $(\mathrm{DR}_{E_{\mathrm{diag}}})$. Let $h_L$ be a proof of $\sfname{RecTExists}$ and $p$ a parameter. If
$\mathrm{pk}(b)$ is true eventually at the bundles of $(\mathcal G_p,\openzeroaLk_p)$, then $\openzeroaNV(D)$ holds
eventually at the bundles of $(\mathcal G_p,\openzeroaLk_p)$.
\formallabel{GlobalTheory.T.puNZ_of_truePack}\ \formallabel{GlobalTheory.MiTop.puNZ_member_alt}
\end{proposition}

\begin{theorem}[The multiplicity identity for true packages]\label{pf:open0a:mitrue}
Assume \ref{pf:open0a:P1}--\ref{pf:open0a:P13}. Let $W$ be a real package over an extra datum of a datum $D$. If $W$ is
true, $r_S\ge1$ and $\openzeroaNV(D)$ holds, then $\openzeroaMI(X_W,e_{V,W},s_{EE,W})$.
\formallabel{GlobalTheory.T.main_theorem_from_published_results_alt_alt_alt_alt_alt}\
\formallabel{GlobalTheory.Fam.miAt_of_truePack_alt_alt}
\end{theorem}

\subsubsection{Proof of the main theorem}
We first prove the main theorem with the conclusion of \cref{pf:open0a:points} as an additional hypothesis.

\begin{proposition}\label{pf:open0a:mainpt}
Under the hypotheses of \cref{pf:open0a:main}, if moreover the point conditions hold eventually at the bundles of
$(\mathcal G_p,\openzeroaLk_p)$, then the conclusion of \cref{pf:open0a:main} holds.
\formallabel{GlobalTheory.T.main_theorem_published_results_final_alt_alt}\
\formallabel{GlobalTheory.Fam.miStar_row_alt_alt}
\end{proposition}

\begin{proof}
Write $\openzeroaEv(E)$ for $\openzeroaEv_{\mathcal G_p,\openzeroaLk_p}(E)$; by hypothesis $\openzeroaEv(\mathcal P)$.

\emph{Step 1: true packages.} We show $\openzeroaEv(b\mapsto\mathrm{pk}(b)\text{ is true})$ by
\cref{pf:open0a:rowEcalc}(b) applied to $\openzeroaEv(\mathcal P)$. Let $K$ be a context, $\beta$ a good point in
context $K$, $c\in\R$ and $b$ a bundle over $(D_\beta,O_\beta,c)$ with $\openzeroaLk_p(K,\beta,c,b)$ and
$\mathcal P(b)$. By \cref{pf:open0a:def:link} there are a weighted point $q$ in context $K$ and a bundle $b'$ over
$(D_q,O_q,c)$ with a bundle witness, such that $(D_\beta,O_\beta,b)=(D_q,O_q,b')$. By this equality $\mathcal P(b')$
holds, and $\mathrm{pk}(b)$ is true if and only if $\mathrm{pk}(b')$ is true. The reading data of $q$ admit a reading
structure by hypothesis, so $\mathrm{pk}(b')$ is true by \cref{pf:open0a:truepack}.

\emph{Step 2: periods.} By Step~1 and \cref{pf:open0a:punz}, $\openzeroaEv(b\mapsto\openzeroaNV(D))$.

\emph{Step 3.} By \cref{pf:open0a:rowEcalc}(a), $\openzeroaEv$ holds for the conjunction of $\mathcal P(b)$, the
truth of $\mathrm{pk}(b)$ and $\openzeroaNV(D)$. At a bundle where the conjunction holds, $\mathrm{pk}(b)$ is true,
$r_S\ge1$ (the last point condition) and $\openzeroaNV(D)$ holds, so \cref{pf:open0a:mitrue} gives
$\openzeroaMI(X_{\mathrm{pk}(b)},e_{V,\mathrm{pk}(b)},s_{EE,\mathrm{pk}(b)})=\openzeroaMI(X_b,e_{V,b},s_{EE,b})$, and
\cref{pf:open0a:rowEcalc}(b) concludes.
\end{proof}

\begin{proof}[Proof of \cref{pf:open0a:main}]
By \cref{pf:open0a:points} (which uses \ref{pf:open0a:P1}--\ref{pf:open0a:P5}), the additional hypothesis of
\cref{pf:open0a:mainpt} holds; apply \cref{pf:open0a:mainpt}.
\end{proof}


%==============================================================================
\subsection{Non-vanishing of periods}\label{sec:pf:open0a:punz}

In this subsection we prove \cref{pf:open0a:punz}. The heart of the argument is a pointwise statement
(\cref{pf:open0a:punzsdiag}): for a true datum satisfying the standing size conditions, the diagonal bound forces some
ramified vector of every word to have a non-zero period. We use the notation of \cref{sec:pf:open0a:words}.

\subsubsection{A kernel bound}

\begin{lemma}\label{pf:open0a:kerHS}
Let $V$ be a finite-dimensional complex inner product space, $K\colon V\to V$ a linear map and $N\subseteq V$ a
subspace with $K(N)=0$. Then $\dim N\le\openzeroaHS{K-\mathrm{id}}^2$.
\formallabel{GlobalTheory.T.finrank_ker_le_hilbertSchmidtNormSq_sub_id}
\end{lemma}

\begin{proof}
Choose an orthonormal basis $e_1,\dots,e_n$ of $V$ such that $e_1,\dots,e_d$ is an orthonormal basis of $N$,
$d=\dim N$. Then $\openzeroaHS{K-\mathrm{id}}^2=\sum_{i=1}^n\|(K-\mathrm{id})e_i\|^2\ge\sum_{i=1}^d\|{-e_i}\|^2=d$.
\end{proof}

\begin{lemma}\label{pf:open0a:kerPu}
For every word $u$ of a datum, $\dim\ker P_u\le\openzeroahs(u)$.
\formallabel{GlobalTheory.Core.Rk.finrank_ker_Pu_le_hsDev}
\end{lemma}

\begin{proof}
$K_u$ is a multiple of $P_u^*P_u$ and vanishes on $\ker P_u$; apply \cref{pf:open0a:kerHS} with $V=C_u$.
\end{proof}

\subsubsection{The ramified share}

\begin{lemma}\label{pf:open0a:ramshare}
If $D$ is true, then $\mathrm{Rs}_u$ holds for every word $u$: $d_U(u)=d_0(u)\,(1-\frac1{T_u-1})^f$.
\formallabel{GlobalTheory.Core.Rk.ramShare_of_trueD}\ \formallabel{GlobalTheory.Core.Rk.ramShare_of_freeT}
\end{lemma}

\begin{proof}
A world over the member construction of $D$ provides level-zero forms for every word $u$, and these give a \emph{free
torus structure} on $M_0(u)$ (\lref{Formalism.L2.GlobalTheoryTheory.LocalLatticeBases.freeT_0}): a finite set
$\mathcal X$ with a free action of the group $\prod_v\F_v^\times$, $v$ running over the places above $T_u$ (each
residue field $\F_v$ has $T_u$ elements), and a linear isomorphism $M_0(u)\cong\C^{\mathcal X}$ which intertwines the
torus operators on $M_0(u)$ with the translations of $\mathcal X$. For such a structure one has the identity of
integers $d_U(u)\,(T_u-1)^f=d_0(u)\,(T_u-2)^f$
(\lref{Formalism.MainNotions.product_first_zeroth_parts.rank_nat}; the ramified subspace corresponds to the
functions on $\mathcal X$ whose sums over the orbits of each factor $\F_v^\times$ vanish, and their dimension is
counted). Since $T_u\ge2$ is prime, $T_u-1>0$, and dividing by $(T_u-1)^f$ gives
$d_U(u)=d_0(u)\bigl(\frac{T_u-2}{T_u-1}\bigr)^f=d_0(u)\bigl(1-\frac1{T_u-1}\bigr)^f$.
\end{proof}

\begin{lemma}\label{pf:open0a:rambig}
If $\mathrm{Rs}_u$, $d_0(u)>0$ and $T_u>2$, then $\mathrm{Rb}_u(f/(T_u-1))$.
\formallabel{GlobalTheory.Core.Rk.ramBig_of_ramShare}
\end{lemma}

\begin{proof}
Put $a=-1/(T_u-1)$; then $-1<a<0$, so $1+a>0$ and $d_U(u)=d_0(u)(1+a)^f>0$. Bernoulli's inequality
$(1+a)^f\ge1+fa$, valid for real $a\ge-2$ and integers $f\ge0$, gives $d_U(u)\ge d_0(u)\bigl(1-f/(T_u-1)\bigr)$.
\end{proof}

\subsubsection{Numerical lemmas}

\begin{lemma}\label{pf:open0a:window}
Let $u$ be a word of a datum $D$.
\begin{enumerate}[label=(\alph*)]
\item If the rank identity holds at $u$, then $d_0(u)\le e\,\Gamma_u^{-2}D_*$.
\item $\log X_u-2\log\Gamma_u/(fk)\le\log T_u$.
\item For all $L,C\in\R$: if the rank identity holds at $u$ and $1+2C\openzeroamfL(L)\le2\log\Gamma_u$, then
$\mathrm{Sm}_u(L,C)$.
\end{enumerate}
\formallabel{GlobalTheory.Core.Rk.levelPrime_window}
\end{lemma}

\begin{proof}
By \eqref{pf:open0a:eq:Tstar}, $\log T_u^*=\log X_u-2\log\Gamma_u/(fk)$, and (b) follows from $T_u^*\le T_u$. Also
$(T_u^*)^{fk}=D_*(u)/(d_u^*\Gamma_u^2)$. Since $1+y\le e^y$, $(1+\frac1{fk})^{fk}\le e$, and the upper end of the
window gives
\[
T_u^{fk}\le\Bigl(1+\frac1{fk}\Bigr)^{fk}(T_u^*)^{fk}\le e\,\frac{D_*(u)}{d_u^*\Gamma_u^2}.
\]
If the rank identity holds, then, as $T_u^k\ge1$,
$d_0(u)=d_u^*(T_u^k-1)^f\le d_u^*T_u^{fk}\le e\,\Gamma_u^{-2}D_*(u)=e\,\Gamma_u^{-2}D_*$; this is (a). Under the
hypotheses of (c), $d_0(u)\le D_*\exp(1-2\log\Gamma_u)\le D_*\exp(-2C\openzeroamfL(L))$.
\end{proof}

\begin{lemma}\label{pf:open0a:diagcor}
Let $u$ be a word of a datum $D$.
\begin{enumerate}[label=(\alph*)]
\item If $A\ge0$, $\mathrm{TT}_u(A,x)$ and $d_0(u)\le\kappa D_*$, then $\mathrm{Df}_u(A\kappa+x)$.
\item If $C\openzeroamfL(L)\ge0$, $\Delta_u(L,C)$ and $\mathrm{Sm}_u(L,C)$, then $\mathrm{Df}_u(\openzeroaeD(L,C))$.
\end{enumerate}
\formallabel{GlobalTheory.Core.Rk.diag_corollary}
\end{lemma}

\begin{proof}
Recall $D_*>0$. (a) From $d_0(u)\le\kappa D_*$ and $d_0(u)\ge0$ we get $d_0(u)^2/D_*\le\kappa d_0(u)$, hence
$\openzeroahs(u)\le A\,d_0(u)^2/D_*+x\,d_0(u)\le(A\kappa+x)\,d_0(u)$.
(b) Put $c=C\openzeroamfL(L)\ge0$. By (a) with $A=4(e^c+1)$, $\kappa=e^{-2c}$ and
$x=4e^{-\openzeroakfk/2}+2e^{-\openzeroakfk}$ we get $\mathrm{Df}_u(e')$ with
\[
e'=4e^{-c}+4e^{-2c}+4e^{-\openzeroakfk/2}+2e^{-\openzeroakfk}\le8e^{-c}+6e^{-\openzeroakfk/2}=\openzeroaeD(L,C),
\]
using $e^{-c}\le1$ and $\openzeroakfk=kf\log k\ge0$. As $d_0(u)\ge0$, the defect bound is monotone in $e$.
\end{proof}

\begin{lemma}\label{pf:open0a:hdnum}
If $\openzeroaSC(D;L,\mathbf C)$ holds, then for every word $u$: $\openzeroamfL(L)\ge5$, $T_u\ge4f+1$, $d_0(u)>0$ and
$fkL\ge5$.
\formallabel{GlobalTheory.Core.Rk.hd_numbers}
\end{lemma}

\begin{proof}
We use \eqref{pf:open0a:eq:HDused}. As $m\ge2$, $f\ge m^4\ge16$ and $k\ge f\ge16$; with $L>1$ this gives
$\openzeroamfL(L)=m^5fL\ge32\cdot16$ and $fkL\ge256$. As $e^L\ge1$, $k^8\le(ke^L)^8\le g_i$ for $i=1,2$; as $f\le k$
and $k\ge4$, $4f\le k^2\le k^8$. Hence $4f\le g_1\le g_1g_2\le N:=(g_1g_2)^2$, and $N\ge2$. Write $T_u=qN+r$ with
$0\le r<N$; the congruence $T_u\equiv1\bmod N$ gives $r=1$, and $q\ne0$ because $T_u\ge2$ is prime. So
$T_u\ge N+1\ge4f+1$. Finally $d_0(u)>0$ is part of \eqref{pf:open0a:eq:HDused}.
\end{proof}

\begin{lemma}\label{pf:open0a:small}
For every datum, $\openzeroakfk/2\ge5$; moreover $e^{-5}<1/128$.
\formallabel{GlobalTheory.Core.Rk.five_le_half_kfk}\ \formallabel{GlobalTheory.Core.Rk.exp_neg_five_lt}
\end{lemma}

\begin{proof}
As $k\ge5>e$, $\log k\ge1$; as $f\ge2$, $kf\ge10$; so $kf\log k\ge10$. Since $e>2.7$, $e^5>2.7^5>128$.
\end{proof}

\begin{lemma}\label{pf:open0a:twoe}
If $C\openzeroamfL(L)\ge5$ and $T_u\ge4f+1$, then $2\,\openzeroaeD(L,C)<1-f/(T_u-1)$.
\formallabel{GlobalTheory.Core.Rk.two_eDiag_lt}
\end{lemma}

\begin{proof}
As $T_u-1\ge4f$, $f/(T_u-1)\le\frac14$. By \cref{pf:open0a:small}, $e^{-C\openzeroamfL(L)}\le e^{-5}$ and
$e^{-\openzeroakfk/2}\le e^{-5}$, so $2\openzeroaeD(L,C)\le28e^{-5}<\frac{28}{128}<\frac34\le1-f/(T_u-1)$.
\end{proof}

\subsubsection{The pointwise statement}

\begin{lemma}\label{pf:open0a:punzat}
Let $u$ be a word of a datum $D$ and $e,\delta\in\R$. If $\mathrm{Df}_u(e)$, $\mathrm{Rb}_u(\delta)$, $e<1-\delta$
and $d_0(u)>0$, then there is $\varphi\in\openzeroaram(u)$ with $P_u(\iota_u^{-1}\varphi)\ne0$.
\formallabel{GlobalTheory.MiTop.puNZ_at}
\end{lemma}

\begin{proof}
Suppose $P_u(\iota_u^{-1}\varphi)=0$ for all $\varphi\in\openzeroaram(u)$. Then $\iota_u^{-1}$ maps
$\openzeroaram(u)$ injectively into $\ker P_u$, so by \cref{pf:open0a:kerPu} and $\mathrm{Df}_u(e)$,
$d_U(u)\le\dim\ker P_u\le\openzeroahs(u)\le e\,d_0(u)$. With $\mathrm{Rb}_u(\delta)$ we get
$(1-\delta)\,d_0(u)\le e\,d_0(u)$, hence $1-\delta\le e$ since $d_0(u)>0$, a contradiction.
\end{proof}

\begin{proposition}\label{pf:open0a:punzsdiag}
Let $D$ be a true datum, $L_2$ size data for $D$ with scale parameter $L=L_2.L$, and $\mathbf C=(C',C_l,C_g,c_b,C_u,
c_T,C_T)$ and $C$ reals with $C\ge1$ and $C+1\le C_\Gamma(L_2)$. Assume $\openzeroaSC(D;L,\mathbf C)$ and, for every
word $u$, the rank identity at $u$ and $\Delta_u(L,C)$. Then $\openzeroaNV(D)$.
\formallabel{GlobalTheory.MiTop.puNZ_of_sDiag}
\end{proposition}

\begin{proof}
Fix a word $u$. By \cref{pf:open0a:hdnum}, $\openzeroamfL(L)\ge5$, $T_u\ge4f+1$ and $d_0(u)>0$. By
\eqref{pf:open0a:eq:Gam}, $\log\Gamma_u=C_\Gamma(L_2)\openzeroamfL(L)$, hence
\begin{gather*}
2\log\Gamma_u\ge2(C+1)\openzeroamfL(L)=2C\openzeroamfL(L)+2\openzeroamfL(L)\ge1+2C\openzeroamfL(L),\\
C\openzeroamfL(L)\ge\openzeroamfL(L)\ge5.
\end{gather*}
By \cref{pf:open0a:window}(c), $\mathrm{Sm}_u(L,C)$, and by \cref{pf:open0a:diagcor}(b),
$\mathrm{Df}_u(\openzeroaeD(L,C))$. Since $f\ge1$, $T_u\ge4f+1>2$, so \cref{pf:open0a:ramshare,pf:open0a:rambig} give
$\mathrm{Rb}_u(f/(T_u-1))$. By \cref{pf:open0a:twoe} and $\openzeroaeD(L,C)\ge0$,
$\openzeroaeD(L,C)\le2\openzeroaeD(L,C)<1-f/(T_u-1)$, and \cref{pf:open0a:punzat} concludes.
\end{proof}

\subsubsection{Proof of the non-vanishing of periods}
Besides $(\mathrm{DR}_{E_{\mathrm{diag}}})$ the proof of \cref{pf:open0a:punz} uses three statements about linked
bundles. In them $L=O.L.L$ is the scale parameter of the size data $O.L$ of the object $O$, and $C_\Gamma$, $C_{YL}$,
$C_{Tu}$ are constants of the system $\kappa.C$ of family constants.

\begin{enumerate}[label=(S\arabic*),leftmargin=*]
\item\label{pf:open0a:S1} (\cref{pf:open0a:sizes}) For all functions
$F,G,H\colon\R^2\to\R$, the following property of $(D,O)$ holds eventually at the bundles of
$(\mathcal G_p,\openzeroaLk_p)$: there are reals $C',C_l,C_g,C_T$ with $F(C_g,C_T)\le C'$, $G(C_g,C_T)\le C_l$,
$H(C_g,C_T)\le L$, $\openzeroaSC(D;L,C',C_l,C_g,\frac14,3,\frac12,C_T)$, $1\le C_\Gamma(L_2)\le C'/8$ for some size data
$L_2$ with scale parameter $L$, the rank identity at every word, the cubic node gap of $D$ (\cref{sec:pf:open0a:points}),
and $r_S\ge1$.
\item\label{pf:open0a:S2} (\cref{pf:open0a:hypL}) For all functions $A,B$ of
systems of family constants, the size inequalities of $D$ at $L$ with the constants
\begin{gather*}
C'=\max(A(\kappa.C),8C_\Gamma),\qquad C_l=\max(B(\kappa.C),40),\qquad C_g=C_{YL}+1,\\
(c_b,C_u,c_T)=(\tfrac14,3,\tfrac12),\qquad C_T=C_{Tu}
\end{gather*}
hold eventually at the bundles of $(\mathcal G_p,\openzeroaLk_p)$.
\item\label{pf:open0a:S3} For every function $J\colon\R^2\to\R$, the property $J(C_{YL}+1,C_{Tu})\le C_\Gamma$ of
$\kappa$ holds eventually in contexts (\cref{pf:open0a:evcgam}).
\end{enumerate}

\begin{lemma}\label{pf:open0a:evcgam}
Statement \ref{pf:open0a:S3} holds.
\formallabel{GlobalTheory.MiTop.evP_cGam}
\end{lemma}

\begin{proof}
The property does not depend on the context, so it suffices to find, for given $\varepsilon,C_8,\dots$, a staged
threshold $K_0$ such that every system $C$ of constants which is large with respect to $K_0$ satisfies
$J(C_{YL}+1,C_{Tu})\le C_\Gamma$. Largeness requires every constant of the list of mid constants of $C$, which
contains $C_\Gamma$, to be at least the value of the mid component of $K_0$ at the list of outer constants of $C$,
which contains $C_{YL}$ and $C_{Tu}$. Take for this mid component the function
$(\text{outer list})\mapsto J(C_{YL}+1,C_{Tu})$ and all other components $0$; take $\varepsilon_0=1$ and all other
bounds $0$.
\end{proof}

\begin{lemma}\label{pf:open0a:trued}
If a real package $W$ over an extra datum of $D$ is true, then $D$ is true, and in the strong sense.
\formallabel{GlobalTheory.Fam.trueD_of_truePack_alt}\ \formallabel{GlobalTheory.MiTop.trueE_of_trueD}
\end{lemma}

\begin{proof}
By definition $D=I.\mathrm{rec}(l,Z)$. A world over $(I,l,Z)$ is assembled from the admissibility of $I$, the pins
and the standard columns contained in the validity of $W$, the spectral ties $W.S$, and the finiteness and identity
conditions contained in the bottom hypotheses. Every world carries the identity condition of its Hida setting, so $D$
is true in the strong sense.
\end{proof}

\begin{proof}[Proof of \cref{pf:open0a:punz}]
Apply $(\mathrm{DR}_{E_{\mathrm{diag}}})$ with $c_b=\frac14$, $C_u=3$, $c_T=\frac12$: for every pair $(C_g,C_T)$ we
obtain reals $C(C_g,C_T)\ge1$, $C'_0(C_g,C_T)$ and $C_{l,0}(C_g,C_T)$ as in \cref{pf:open0a:def:rowdiag}. By
\cref{pf:open0a:rowEcalc}(a) the conjunction of the following properties holds eventually at the bundles of
$(\mathcal G_p,\openzeroaLk_p)$: $\mathrm{pk}(b)$ is true (hypothesis); \ref{pf:open0a:S1} with $F=G=H=0$;
\ref{pf:open0a:S2} with $A(C)=C'_0(C_{YL}+1,C_{Tu})$ and $B(C)=C_{l,0}(C_{YL}+1,C_{Tu})$; and \ref{pf:open0a:S3} with
$J=C+1$ (a property which holds eventually in contexts holds in particular eventually at the bundles). Let $K$ be a
context of type $\kappa$, $\beta$ a good point in context $K$ and $b$ a linked bundle over $(D,O)=(D_\beta,O_\beta)$
at which the conjunction holds. Put $L=O.L.L$, $C_g=C_{YL}+1$, $C_T=C_{Tu}$, $C=C(C_g,C_T)$,
$C'=\max(C'_0(C_g,C_T),8C_\Gamma)$ and $C_l=\max(C_{l,0}(C_g,C_T),40)$. Then:
\begin{itemize}
\item $D$ is true in the strong sense by \cref{pf:open0a:trued};
\item the size data $O.L$ of a point in a context of type $\kappa$ have the constant $C_\Gamma(O.L)=C_\Gamma$ of
$\kappa.C$, so $C+1\le C_\Gamma(O.L)$ by \ref{pf:open0a:S3};
\item $\openzeroaSC(D;L,C',C_l,C_g,\frac14,3,\frac12,C_T)$ holds: its size inequalities are \ref{pf:open0a:S2}, and its
other four components (ties, mid-volume formula, interlacing and congruence conditions, trace frame) do not depend on
the constants and are given by \ref{pf:open0a:S1};
\item $E_{\mathrm{diag}}(D,L,C')$ holds: the rank identity at every word by \ref{pf:open0a:S1}, and $O.L$ has scale
parameter $L$, $8C_\Gamma(O.L)\le C'$ and $C_\Gamma(O.L)\ge C+1\ge1$.
\end{itemize}
Since $C'\ge C'_0(C_g,C_T)$ and $C_l\ge C_{l,0}(C_g,C_T)$, $(\mathrm{DR}_{E_{\mathrm{diag}}})$ gives $\Delta_u(L,C)$ for
every word $u$, and \cref{pf:open0a:punzsdiag} with $L_2=O.L$ gives $\openzeroaNV(D)$. \Cref{pf:open0a:rowEcalc}(b)
concludes.
\end{proof}


%==============================================================================
\subsection{Point conditions and true packages}\label{sec:pf:open0a:points}

The bundle witness of a linked bundle already contains most of the data from which the point conditions and the truth
of the real package are read off; the remaining conditions are size conditions which hold eventually. For a datum $D$,
let $b_1(\tau)<b_2(\tau)<\dots$ be the \emph{nodes} of the Harish-Chandra parameter at the infinite place $\tau$, and
$N_m=\mathrm{rank}(X')+1$. The \emph{cubic node gap} of $D$ is the condition that consecutive nodes differ by at least
$4N_m^3+2t+6$; the \emph{node gap condition} is that there is $N\ge N_m$ such that consecutive nodes differ by at
least $t+2N^3+1$.

\begin{proof}[Proof of \cref{pf:open0a:points}]
We use the following statements.
\begin{enumerate}[label=(\alph*)]
\item Statement \ref{pf:open0a:S1} (\cref{pf:open0a:sizes}).
\item (\lref{GlobalTheory.Core.Rk.HdNc.hd_nc}) There is $C^{\mathrm{nc}}_{l,0}$ such that whenever
$C_l\ge C^{\mathrm{nc}}_{l,0}$ and $\openzeroaSC(D;L,C',C_l,\dots)$ holds, every Hida setting of $D$ satisfies the gap
condition with constant $100\,r^4$. (The formal proof takes $C^{\mathrm{nc}}_{l,0}=1801$ and derives the gap from the
inequalities between consecutive entries of the parameters contained in the size inequalities.)
\item (\lref{GlobalTheory.Fam.PtRowsAt.gapHyp_of_hd}) If $C_l\ge1801$ and $\openzeroaSC(D;L,C',C_l,\dots)$, then
the gap hypothesis of the Poincar\'e amplitude bounds holds for $D$; and
(\lref{GlobalTheory.Core.Rk.finX_of_gap}) under this hypothesis the spectral data of every system of spectral
ties over $D$ satisfy the finiteness condition (by the decay of Wu functions in large gap and the finite
dimensionality of spaces of automorphic forms).
\item (\lref{GlobalTheory.Core.Rk.GapOfNodes.harishChandraNodesFarApart_of_cubic_node_gap}) The cubic node gap
implies the node gap condition.
\item For a weighted point $q$ and a bundle $b'$ with a bundle witness: under \ref{pf:open0a:P1}--\ref{pf:open0a:P5},
the realisation of $b'$ extends and a join of it with $b'.R$ and $b'.S$ exists
(\lref{GlobalTheory.Fam.PtRowsAt.seeEx_and_joinSI_nonempty_of_bWit}); and, without further hypotheses, the
decomposition point condition of $b'.P$, the existence of frames, the existence condition of the real datum, the
primitivity rows and the separation of types hold; these are read off the witness
(\lref{GlobalTheory.Fam.PtRowsAt.dcPtG_of_bWit}, \lref{GlobalTheory.Fam.PtRowsAt.framesEx_of_bWit},
\lref{GlobalTheory.Fam.PtRowsAt.exAt_of_bWit},
\lref{GlobalTheory.Fam.PtRowsAt.prim_links_and_typesSep3_of_bWit}).
\end{enumerate}
Apply (a) with $F=H=0$ and $G\equiv\max(1801,C^{\mathrm{nc}}_{l,0})$, and use \cref{pf:open0a:rowEcalc}(b). At a bundle
$b$ linked to a good point $\beta$ at which the property of (a) holds, choose $q,b'$ and a bundle witness with
$(D_\beta,O_\beta,b)=(D_q,O_q,b')$ (\cref{pf:open0a:def:link}). The first six point conditions hold for $b'$ by (e),
hence for $b$. The constants of (a) satisfy $C_l\ge1801$ and $C_l\ge C^{\mathrm{nc}}_{l,0}$, so (c) gives the
finiteness condition for the spectral data of $b.S$ and (b) gives the gap condition for the Hida setting of $b.S$;
(d) gives the node gap condition from the cubic node gap; and $r_S\ge1$ is part of (a).
\end{proof}

\begin{proof}[Proof of \cref{pf:open0a:truepack}]
We use four statements about the bundle witness.
\begin{enumerate}[label=(\alph*)]
\item (\lref{GlobalTheory.Fam.colsStd_of_bWit_alt}) The construction of $q$ is admissible, pins exist, the
construction is twisted, standard columns exist for the columns of $b.R$, and $\mathrm{pk}(b)$ satisfies the standard
twist condition.
\item (\lref{GlobalTheory.T.characters_match_bundle_alt}) The characters of the package match those of the
construction: the characters $\mu_0$ at the levels $j+1$, $j<m$; the characters $\xi_t$ for every tuple $t$; and the
prime sets $Q_t$ (they are the images, under the selection map of the construction, of the prime set attached to the
character of $t$). Moreover $\mathrm{pk}(b)$ satisfies the level condition, a recursion of the realisation exists,
the realisation is normalised, the row-side family satisfies its admissibility condition, the Hida setting of $b.S$
satisfies the identity condition, and the column dictionaries satisfy the level matching for every word.
\item (\cref{pf:open0a:readat}) If the reading data of $q$ admit a reading structure, then for every word $u$ the
leading read of $\mathrm{pk}(b)$ holds at the identity for some top map.
\item (\lref{GlobalTheory.Fam.pkHyp_of_real_alt}) A real package over a datum of the member construction
satisfies the package hypotheses if: a join exists; the level matching holds; the decomposition point condition holds;
frames exist; the realisation is normalised, extends and admits a recursion; the existence condition of the real
datum holds; the primitivity rows hold and the types are separated; the row-side family is admissible; the identity
condition and the gap condition hold for the Hida setting; the finiteness condition and the node gap condition hold;
and the package is valid and satisfies the geometric conditions.
\end{enumerate}
By (a), (b) and (c), $\mathrm{pk}(b)$ is valid and satisfies the geometric conditions. All hypotheses of (d) are now
available (those not yet mentioned are point conditions of $b$), so $\mathrm{pk}(b)$ satisfies the package hypotheses,
and with (a) and (b) all components of $\mathcal A(\mathrm{pk}(b))$ hold. Finally $D_q$ is the datum
$I.\mathrm{rec}(l,Z)$ of the member construction of $q$, so $\mathrm{pk}(b)$ is true.
\end{proof}


%==============================================================================
\subsection{The multiplicity identity for true packages}\label{sec:pf:open0a:mi}

In this subsection we prove \cref{pf:open0a:mitrue}. Throughout, $W$ is a real package over an extra datum of a datum
$D$, $X=X_W$ its coefficient complex with face $(n_V,\mathrm{Mem},a,X_c,\mathrm{Pt},H,P)$, $e_V=e_{V,W}$,
$s_{EE}=s_{EE,W}$, and $\mathcal T$ the set of tuples of the row-side family $W.P$, with base tuple $\mu_0$.

\subsubsection{Reduction}

\begin{proposition}\label{pf:open0a:allTrue}
Assume \ref{pf:open0a:P1}--\ref{pf:open0a:P13}. Let $D=I.\mathrm{rec}(l,Z)$ be the datum of a member construction and
$W$ a real package over an extra datum of $D$ with $\mathcal A(W)$, $r_S\ge1$ and $\openzeroaNV(D)$. Then
$\openzeroaMI(X_W,e_{V,W},s_{EE,W})$.
\formallabel{GlobalTheory.T.main_theorem_from_published_results_alt_alt_alt_alt}\
\formallabel{GlobalTheory.Fam.miStar2_of_allTrue_alt_alt}
\end{proposition}

\begin{proof}
By \cref{pf:open0a:onto}, the surjectivity along chains $\mathrm{Onto}(D)$ holds. By \cref{pf:open0a:famall} (which
uses \ref{pf:open0a:P2}, \ref{pf:open0a:P3}, \ref{pf:open0a:P4}, \ref{pf:open0a:P13}), the family conditions
$\mathrm{FC}(W)$ hold. By \cref{pf:open0a:served}, the served condition $\mathrm{SV}(W)$ holds. The package hypotheses
contain the bottom hypotheses and the level hypothesis, and \cref{pf:open0a:mistarpunz} (which uses \ref{pf:open0a:P2},
\ref{pf:open0a:P4}, \ref{pf:open0a:P12}) concludes.
\end{proof}

\begin{proof}[Proof of \cref{pf:open0a:mitrue}]
If $W$ is true, then $D=I.\mathrm{rec}(l,Z)$ for a member construction and $\mathcal A(W)$ holds; apply
\cref{pf:open0a:allTrue}.
\end{proof}

\subsubsection{Surjectivity along chains}
A datum $D$ has \emph{chains} $c$ (the iterated theta lifts of \eqref{eq:families:chain}) and, at each level $k\le m$
of a chain, \emph{localisations} $s$, each with a space $s.\mathrm{CC}(c)$ of cusp forms; a \emph{step} from a
localisation at level $k$ to one at level $k+1$ carries a theta map back to level $k$, and a theta map
$\Theta\colon s.\mathrm{CC}(c)\to s_0.\mathrm{CC}(c)$ is \emph{reachable in $k$ steps} from a localisation $s_0$ at
level $0$ if it is the composite of the theta maps of $k$ admissible steps starting at $s_0$.

\begin{definition}\label{pf:open0a:def:onto}
$\mathrm{Onto}(D)$: if the node gap condition holds for $D$, then for every admissible chain $c$, every localisation
$s_0$ at level $0$, every $k\le m$ and every theta map $\Theta$ reachable in $k$ steps from $s_0$, $\Theta$ is
surjective.
\formal{GlobalTheory.ChainOnto.OntoAll}
\end{definition}

\begin{proposition}\label{pf:open0a:onto}
Assume \ref{pf:open0a:P1}--\ref{pf:open0a:P11}. Then $\mathrm{Onto}(D)$ holds for every datum $D$.
\formallabel{GlobalTheory.T.ontoAll_of_theta_doubling}\ \formallabel{GlobalTheory.ChainOnto.ontoAll_of_hecke_alt}
\end{proposition}

\begin{proof}
The proof combines three statements.
\begin{enumerate}[label=(\alph*)]
\item (\lref{GlobalTheory.T.heckeForm_of_doubling}) Under \ref{pf:open0a:P2} and
\ref{pf:open0a:P7}--\ref{pf:open0a:P11}, every datum satisfies a Rallis inner product formula for the lifts of the
chain: for a Weil representation of a pair $(A,B)$ in the convergent range, Fock vectors $\Phi_1,\Phi_2$ unramified
outside a finite set $S_f$ of places, and cusp forms $u_1,u_2$ on $\Uni(A)$ ($u_1$ with Hecke parameter
$\mathrm{par}$, $u_2$ spherical outside $S_f$) for which the doubling integrand converges, the lifts converge and
\[
\bigl\langle\theta(\Phi_1,u_1),\theta(\Phi_2,u_2)\bigr\rangle_{\Uni(B)}
=\mathrm{vol}\cdot\frac{l}{d^{S_f}(t/2)}\int(\text{Weil coefficient})\cdot(\text{matrix coefficient}),
\]
the integral being over the complement of $S_f$, where $l$ is the value at $t/2$ of the partial $L$-function
$L^{S_f}(s+\frac12,\mathrm{par}\otimes\chi_B)$ and $d^{S_f}$ is the normalising factor of the doubling integral.
\item (\lref{GlobalTheory.T.thetaChainUpwardLift_heckeParam_transferred}) Under \ref{pf:open0a:P2} and
\ref{pf:open0a:P6}, the level facts of the chain (\lref{GlobalTheory.T.thetaChain_level_properties}: the level
groups and level characters of the localisations are compatible with the local stabilisers) imply the
\emph{parameter transfer}: if a vector of a localisation at level $k$ has Hecke parameter $B$ and its upward lift is
non-zero, then the lift has the parameter obtained from $B$ by the step (adjoining a segment of length $t$).
\item (\cref{pf:open0a:ontoclassical}) Under \ref{pf:open0a:P1}--\ref{pf:open0a:P5}, the inner product formula of (a)
and the parameter transfer of (b) imply $\mathrm{Onto}(D)$.
\end{enumerate}
Applying (c) to the outputs of (a) and (b) proves the proposition.
\end{proof}

\subsubsection{Rows, columns, bottom data and chain data}\label{sec:pf:open0a:objects}
We now describe the objects in terms of which the family, joint and final conditions are formulated. Let $G$ be the
group of adelic points of $\Uni(Y_0)$, $Y_0$ the hermitian space of level $0$ of $D$, $G_f\subseteq G$ the subgroup of
finite adelic points and $Z(G)$ the centre of $G$. The realisation $W.Q$ provides an element $t_S\in G$, an archimedean
family $\theta$, discrete series models $M$ at the infinite places with the space $M^{\infty}$ of smooth tensors, a
level set $J\subseteq G$ and a level character $\chi\colon G\to\C$ at the true level, and an additive character $\psi$.
The spectral ties $W.S$ provide a Hida setting with $L^2$-space $L^2_X$, a closed subspace $H_X\subseteq L^2_X$ of cusp
forms, an isomorphism $\iota$ of its group with $G$, and for every word $u$ a column space $\mathrm{Cu}(u)$ with period
functions $\varpi_C\in L^2_X$ for $C\in\mathrm{Cu}(u)$.
\begin{itemize}
\item The \emph{base row space} $\openzeroaRow$ is the space of multiplicity maps of type $(\theta,M,J,\chi)$: linear
maps from $M^{\infty}$ to functions on $G$ with the equivariance properties of that type.
\item The realisation provides a finite-dimensional Hilbert space $H_E$ with an injective linear map $e_E$ from $H_E$ to
the multiplicity maps of $M$; for $e\in H_E$, $\mathrm{rf}(e)\colon M^{\infty}\to\C^G$ is the induced linear map,
the \emph{row form} of $e$. The space $E$ of the coefficient complex $X$ is $H_E$, with scalar multiplication
conjugated if the twist bit $s\in\{0,1\}$ of $W.P$ is $1$; we write $\openzeroacj(z)=\bar z$ if $s=1$ and
$\openzeroacj(z)=z$ if $s=0$. Then $\iota_0=\openzeroacj\circ\iota_{\mathit{LL}}$, where $\iota_{\mathit{LL}}$ is the
given embedding of $\mathit{LL}$. The true part $\openzeroaTrue$ of $E$ consists of the $e$ such that $e_E(e)$ lies in
the good part of the multiplicity maps at the true level.
\item For a tuple $t$ and a row $F\in\mathit{LL}^{n_V(t)}$, $\mathrm{rf}(t,F)$ is the row form of the theta image
$\Theta_t(j_R(t)(e_V(t)F))\in H_E$ of the row.
\item For $g\in G$, $f\colon G\to\C$, $C\in\mathrm{Cu}(u)$ and $b\in\C$, the \emph{pairing relation}
$\mathrm{Pair}_u(g,f,C,b)$ means: there is $\phi\in L^2_X$ equal almost everywhere to the translate
$x\mapsto f(\iota(x)g)$, and $b=\langle\bar\phi,\varpi_C\rangle$.
\item Let $M_u$ be the column space of the word $u$; it is identified with the ramified subspace by a linear isomorphism
$e_u\colon M_u\to\openzeroaram(u)$. $C\in\mathrm{Cu}(u)$ is a \emph{column form} of $\varphi\in M_u$
($\mathrm{CF}_u(C,\varphi)$) if the column dictionary of $u$ maps $C$ to $e_u(\varphi)$.
\item For $\varphi\in M_u$, a tuple $t$ and a row $F$ of $t$, $\mathrm{pc}_t(\varphi,F)\in\C$ is the column pairing
of the datum between $\varphi$ and the row $F$, read in the ambient of the tuple $t$; it is linear in $\varphi$.
\end{itemize}

\begin{definition}[Bottom data]\label{pf:open0a:def:bottom}
A \emph{bottom datum} $d$ for $(W,u)$ consists of a Bessel datum $d.B$ for the Hida setting, the word $u$, $\psi$ and
the space $X_0$ of level $0$, elements $d.b,d.k\in G$ and $d.a_0,d.a\in\Uni(X_0)(\A)$. Its \emph{periods} are
\[
d.\mathrm{per}(\varphi,T)(y,v)=\mathrm{Bes}_{d.B}\bigl(\mathrm{rd}_0(e_u\varphi)(y),\ R_{d.b}(T(v))\bigr)
\]
for $\varphi\in M_u$, linear maps $T\colon M^{\infty}\to\C^G$, reading points $y$ and $v\in M^{\infty}$: the
Bessel period of the level of $d.B$ (with character $\psi$) between the reading at level $0$ of $e_u\varphi$ at $y$
and the right translate by $d.b$ of $T(v)$. For $z\in G$, $d$ is \emph{pinned at} $z$ if $d.B$ is adapted to the
transport weights of $u$, $d.a_0$ is an adelic transport correction, $d.a$ is a finite central element of
$\Uni(X_0)(\A)$, $d.k\in J$, and $\mathrm{emb}_{d.B}(d.a_0\,d.a)\,d.b=\iota(1,x_u)\,z\,t_S\,d.k$ in $G$, where
$(1,x_u)$ is the element of the group of the Hida setting attached to $u$.
\formal{GlobalTheory.MiTop.BottomData}\ \formal{GlobalTheory.MiTop.BottomData.Pin}
\end{definition}

\begin{definition}[Relations]\label{pf:open0a:def:rel}
Let $d$ be a bottom datum for $(W,u)$, $z\in G$, $U\colon M_u\to M_u$ linear and $c\colon\mathcal T\to\C$.
\begin{enumerate}[label=(\alph*)]
\item The \emph{bottom relation} $\mathrm{BR}(z,d)$: there is a linear isometry $R$ of $M_u$ such that for all
$\varphi,\varphi'\in M_u$, $T,T'\in\openzeroaRow$ and $\alpha,\alpha'\in\C$ with
$\alpha\,d.\mathrm{per}(\varphi,T)(y,v)+\alpha'\,d.\mathrm{per}(\varphi',T')(y,v)=0$ for all $y,v$, and for all
$v\in M^{\infty}$, column forms $C$ of $R\varphi$ and $C'$ of $R\varphi'$, and $b,b'$ with
$\mathrm{Pair}_u(zt_S,T(v),C,b)$ and $\mathrm{Pair}_u(zt_S,T'(v),C',b')$, one has $\alpha b+\alpha'b'=0$.
\item The relation $\mathrm{Rel}_d(U,c)$: for every $T\in\openzeroaRow$ there is $p\colon M_u\to\C$ such that for all
$\varphi_0,\varphi\in M_u$, tuples $t$, rows $F$ of $t$ and $\alpha,\alpha'\in\C$ with
$\alpha\,p(\varphi_0)+\alpha'\,c(t)\,\mathrm{pc}_t(\varphi,F)=0$ one has
$\alpha\,d.\mathrm{per}(U\varphi_0,T)(y,v)+\alpha'\,d.\mathrm{per}(U\varphi,\mathrm{rf}(t,F))(y,v)=0$ for all $y,v$.
\item $c$ is a system of \emph{unit coefficients}, $\mathrm{UC}(c)$, if $c(t)\ne0$ for all $t$ and, for every finite
place $w$ of $\mathit{LL}$ and every $k\in\mathit{LL}$ with $\iota_{\mathit{LL}}(k)=c(\mu(w))/c(\mu_0)$, $v_w(k)=0$.
\item The \emph{joint condition} $\mathrm{JC}(W)$: there is $z\in G_f\cap Z(G)$ such that for every word $u$ there are a
bottom datum $d$ pinned at $z$, a linear map $U$ of $M_u$, $r>0$ with $\|U\varphi\|=r\|\varphi\|$ for all $\varphi$, and
$c\colon\mathcal T\to\C$ with $\mathrm{UC}(c)$ and $\mathrm{Rel}_d(U,c)$.
\end{enumerate}
\formal{GlobalTheory.MiTop.BottomRelation}\ \formal{GlobalTheory.MiTop.JointCondition}
\end{definition}

\paragraph{Family characters and chain data.}
For a word $u$, a \emph{family character} $f$ for $(W,u)$ consists of a geometric chain datum $f.g$, with base points
$f.g.b(i)\in G$, a top map, and choices of rows and columns for every tuple $t$. It determines a commuting chain $f.c$,
with Bessel data $f.c.B(0,n)$ at the bottom, and for every tuple $t$ a \emph{tuple chain} $f_t$. Since $m$ is even we
put $N=m/2$ and $n=N-1$. We use the following components of $f_t$:
\begin{itemize}
\item complex vector spaces $\mathrm{Col}_0(f_t)\supseteq\mathrm{Ram}_0(f_t)$ (columns and ramified columns at level
$0$) and $\mathrm{Col}_m(f_t)$ (columns at level $m$), a set $\mathrm{CT}(f_t)$ of \emph{top forms}, and the embedding
$\mathrm c_0^t\colon M_u\to\mathrm{Col}_0(f_t)$;
\item linear maps $\mathrm{up}^t,\mathrm{col}^t\colon\mathrm{Col}_0(f_t)\to\mathrm{Col}_m(f_t)$ (the lift along the chain
and the transport of columns) and the \emph{top pairing}
$\mathrm{tail}^t\colon\mathrm{Col}_m(f_t)\times\mathrm{CT}(f_t)\to\C$;
\item for $G\in\mathrm{CT}(f_t)$, its \emph{descent} $\mathrm{desc}^t(G)$ to level $0$, a linear map
$M^{\infty}\to\C^G$, and the \emph{bottom periods} $\mathrm{Per}^t(\psi,G)(y,v)\in\C$ of $\psi\in\mathrm{Col}_0(f_t)$
against the descent of $G$;
\item constants $c_{\mathrm{top}}^t$ (at level $N$) and $c_{\mathrm{bot}}^t$; the latter is a product of constants which are
all non-zero when $f_t$ is admissible (satisfies its structure conditions).
\end{itemize}
The \emph{family laws} $\mathrm{Laws}(f)$ are the conjunction of: $f.c$ and all $f_t$ are admissible, the $f_t$ are
locally finite, and the reality, valuation, last-step, kernel, period and reading laws hold for the family
$(f_t)_{t\in\mathcal T}$.

\begin{definition}[Link conditions]\label{pf:open0a:def:links}
Let $f$ be a family character for $(W,u)$ and $d$ a bottom datum. Below, $t$ runs over tuples, $\varphi$ over $M_u$, $F$
over rows of $t$, $G$ over $\mathrm{CT}(f_t)$, and $\mathrm{lead}(t,F)$ is the leading linear map of the row $F$, a
possible leading component of top forms.
\begin{enumerate}[label=(L\arabic*),leftmargin=*]
\item\label{pf:open0a:L1} \emph{Bottom compatibility}: $\mathrm c_0^t(\varphi)\in\mathrm{Ram}_0(f_t)$ and
$\mathrm{Per}^t(\mathrm c_0^t\varphi,G)(y,v)=d.\mathrm{per}(\varphi,\mathrm{desc}^t(G))(y,v)$ for all $y,v$.
\item\label{pf:open0a:L2} \emph{Tail identity}: if $G$ has leading component $\mathrm{lead}(t,F)$, then
$\mathrm{tail}^t(\mathrm{col}^t(\mathrm c_0^t\varphi),G)=\mathrm{pc}_t(\varphi,F)$.
\item\label{pf:open0a:L3} \emph{Descent}: for all $t,F$ there is $G$ with leading component $\mathrm{lead}(t,F)$ and
$\mathrm{desc}^t(G)=\mathrm{rf}(t,F)$.
\item\label{pf:open0a:L4} \emph{Surjectivity}: for every $T\in\openzeroaRow$ there is $G_0\in\mathrm{CT}(f_{\mu_0})$ with
$\mathrm{desc}^{\mu_0}(G_0)=T$.
\item\label{pf:open0a:L5} \emph{Commutation}: there are a linear automorphism $S$ of $M_u$, $r>0$ with
$\|S\varphi\|=r\|\varphi\|$ for all $\varphi$, and $s_0\in\C^\times$ such that
$\mathrm{tail}^t(\mathrm{up}^t(\mathrm c_0^t\varphi),G)=s_0\,\mathrm{tail}^t(\mathrm{col}^t(\mathrm c_0^t(S\varphi)),G)$
for all $t,\varphi,G$.
\item\label{pf:open0a:L6} \emph{Ratio condition}: for every $t$, $c_{\mathrm{top}}^t\ne0$ and
$\mathrm{UR}_{\mathcal Q_t}\bigl(c_{\mathrm{top}}^t/c_{\mathrm{bot}}^t,\ c_{\mathrm{top}}^{\mu_0}/c_{\mathrm{bot}}^{\mu_0}\bigr)$,
where $\mathcal Q_t$ is the set of rational primes contained in some prime of the set $Q_t$ of primes of $E$, and
$\mathrm{UR}_{\mathcal S}(x,y)$ (``$x/y$ is a root of an $\mathcal S$-unit'') means: $y\ne0$, and there are $n_0\ge1$
and $\rho\in\Q^\times$ with $(x/y)^{n_0}=\rho$ such that every prime $p$ with $v_p(\rho)\ne0$ lies in $\mathcal S$.
\end{enumerate}
The \emph{family conditions} $\mathrm{FC}(W)$: there is $z\in G_f\cap Z(G)$ such that for every word $u$ there are a
family character $f$ and a bottom datum $d$ pinned at $z$ with $\mathrm{Laws}(f)$ and
\ref{pf:open0a:L1}--\ref{pf:open0a:L6}. The \emph{served condition} $\mathrm{SV}(W)$: for every finite place $w$ of
$\mathit{LL}$ and every $p\in\mathcal Q_{\mu(w)}$, $w\nmid p$.
\formal{GlobalTheory.MiTop.FamilyAllConditions}\ \formal{GlobalTheory.MiTop.FamilyServedCondition}
\end{definition}

\subsubsection{The family conditions}

\begin{proposition}\label{pf:open0a:famall}
Assume \ref{pf:open0a:P2}, \ref{pf:open0a:P3}, \ref{pf:open0a:P4} and \ref{pf:open0a:P13}. Let $W$ be a real package
over an extra datum of $D=I.\mathrm{rec}(l,Z)$ with $\mathcal A(W)$, $r_S\ge1$ and $\mathrm{Onto}(D)$. Then
$\mathrm{FC}(W)$.
\formallabel{GlobalTheory.T.true_geometry_family_row_data_alt}\ \formallabel{GlobalTheory.MiTop.famAllQ_of_allTrue_alt}
\end{proposition}

The proof assembles the following results; those proved in \cref{sec:pf:open0a:deferred} make up most of the formal
proof of the main theorem. A \emph{geometric family} $a$ for $(W,u)$ consists of a
geometric chain datum $a.g$, a top map, a non-archimedean local field with a type datum, a commuting vector datum, local
group data indexed by the type places of $W$, local geometric data at the primes of the sets $Q_t$, a bottom datum
$a.d$ and constants; together with row and column choices it determines a family character
$a.\mathrm{fam}(\mathrm{row},\mathrm{col})$. A Bessel datum is \emph{framed} if its vectors are, up to a non-zero
scalar and powers of $\pm g$ ($g$ the type prime of the Hida setting), the standard vectors of the Hida setting.

\begin{enumerate}[label=(F\arabic*),leftmargin=*]
\item\label{pf:open0a:F1} (\cref{pf:open0a:botdat}) A central pin: there is $z\in G_f\cap Z(G)$ such that placed framed
Bessel data extend to bottom data pinned at $z$.
\item\label{pf:open0a:F2} (\cref{pf:open0a:geoall}) For every word there is an admissible and true geometric family with
identity transfer matrix and framed bottom Bessel datum.
\item\label{pf:open0a:F3} (\lref{GlobalTheory.T.exists_pairwise_same_choices_at_geo}) Under \ref{pf:open0a:P2}
and the package hypotheses: for every geometric chain datum there are row choices $\mathrm{row}(t)$,
$t\in\mathcal T$, all admissible and pairwise of the same shape.
\item\label{pf:open0a:F4} (\cref{pf:open0a:colside}) For every system of family constants there are column choices
satisfying the column condition and the column output condition.
\item\label{pf:open0a:F5} (\cref{pf:open0a:lawslinks}) If \ref{pf:open0a:F4} holds for every system of constants, some
column choice gives a family character with its row conditions, the family laws, the chart condition at level $0$,
\ref{pf:open0a:L2}, \ref{pf:open0a:L5} and \ref{pf:open0a:L6}.
\item\label{pf:open0a:F6} (\cref{pf:open0a:isfam}) The row conditions imply that the commuting chain is admissible.
\item\label{pf:open0a:F7} (\cref{pf:open0a:placed}) For an admissible true geometric family the placement conditions
hold for $(a.\mathrm{fam}(\mathrm{row},\mathrm{col}),a.d.B)$ whenever its commuting chain is admissible.
\item\label{pf:open0a:F8} (\lref{GlobalTheory.T.linkZero_of_chartZero'}) If $f$ satisfies the chart condition at
level $0$, the level of $d.B$ equals $f.c.B(0,r_S)$ whenever $2r_S+2=m$, and $d.b=f.g.b(0)$, then
\ref{pf:open0a:L1} holds.
\item\label{pf:open0a:F9} (\lref{GlobalTheory.T.linkDown_of_famCh}) Under the package hypotheses, the row conditions
of $f$ imply \ref{pf:open0a:L3}.
\item\label{pf:open0a:F10} (\cref{pf:open0a:linkonto}) Under the package hypotheses and $\mathrm{Onto}(D)$, the row
conditions of $f$ imply \ref{pf:open0a:L4}.
\end{enumerate}

\begin{proof}[Proof of \cref{pf:open0a:famall}]
Take $z$ from \ref{pf:open0a:F1} (the package hypotheses are part of $\mathcal A(W)$). Let $u$ be a word. Choose $a$ by
\ref{pf:open0a:F2}, row choices for $a.g$ by \ref{pf:open0a:F3}, and apply \ref{pf:open0a:F4} for every system $k$ of
family constants; then \ref{pf:open0a:F5} gives column choices, and we put $f=a.\mathrm{fam}(\mathrm{row},\mathrm{col})$.
By \ref{pf:open0a:F6}, $f.c$ is admissible, so \ref{pf:open0a:F7} gives the placement conditions for $(f,a.d.B)$; as
$a.d.B$ is framed, \ref{pf:open0a:F1} gives a bottom datum $d$ pinned at $z$ with $d.B=a.d.B$ and $d.b=f.g.b(0)$. The
bottom clause of the admissibility of $a$ says that the level of $a.d.B$ equals $f.c.B(0,r_S)$ when $2r_S+2=m$; hence
\ref{pf:open0a:F8} gives \ref{pf:open0a:L1}. \ref{pf:open0a:F9} and \ref{pf:open0a:F10} give \ref{pf:open0a:L3} and
\ref{pf:open0a:L4}, and \ref{pf:open0a:F5} gives $\mathrm{Laws}(f)$, \ref{pf:open0a:L2}, \ref{pf:open0a:L5} and
\ref{pf:open0a:L6}.
\end{proof}

\subsubsection{The served condition}

\begin{lemma}\label{pf:open0a:served}
Every real package $W$ satisfies the served condition $\mathrm{SV}(W)$. (The formal statement also assumes the bottom
hypotheses and the existence of a recursion of the realisation, which are not used.)
\formallabel{GlobalTheory.MiTop.servedQ_of_bot_alt}\ \formallabel{GlobalTheory.MiTop.served'}
\end{lemma}

\begin{proof}
We use the following properties of the character data of the realisation $W.Q$, which are part of its definition.
For each tuple $t$ there is a Hecke character $\xi_t$ of $E$, with $\xi_{\mu_0}=1$, which is ramified at every prime of
$Q_t$. Each finite place $w$ of $\mathit{LL}$ has a residue characteristic $\ell(w)$. If $w$ is not twisted, then
$\mu(w)=\mu_0$. If $w$ is twisted, then $\xi_{\mu(w)}=\delta_{\ell(w)}\psi_w$, where $\psi_w$ is unramified at every
prime of $E$ above $\ell(w)$, and $\delta_\ell=\delta'$ if $\ell=p_{\mathrm{de}}$ and $\delta_\ell=\delta$ otherwise;
here $\delta$ (resp.\ $\delta'$) is unramified at every prime not above $p_{\mathrm{de}}$ (resp.\ $p'_{\mathrm{de}}$),
and $p_{\mathrm{de}}\ne p'_{\mathrm{de}}$ are auxiliary primes of the host.

Let $w$ be a finite place of $\mathit{LL}$, $p$ a prime and $\mathfrak P\in Q_{\mu(w)}$ with $p\in\mathfrak P$, and
suppose $w\mid p$. Then $\ell(w)=p$, and $p$ is the only rational prime contained in $\mathfrak P$ (two distinct primes
are coprime). The character $\xi_{\mu(w)}$ is ramified at $\mathfrak P$. If $w$ is not twisted,
$\xi_{\mu(w)}=\xi_{\mu_0}=1$, a contradiction. If $w$ is twisted, $\psi_w$ is unramified at $\mathfrak P$; and
$\delta_p$ is unramified at $\mathfrak P$: if $p=p_{\mathrm{de}}$, then $\delta_p=\delta'$ and $\mathfrak P$ lies over
$p\ne p'_{\mathrm{de}}$; otherwise $\delta_p=\delta$ and $\mathfrak P$ lies over $p\ne p_{\mathrm{de}}$. So
$\xi_{\mu(w)}=\delta_p\psi_w$ is unramified at $\mathfrak P$, again a contradiction.
\end{proof}

\subsubsection{Seesaw transport relations and the joint condition}
The family laws imply the \emph{seesaw transport relations}, which transfer linear relations between top pairings to
linear relations between bottom periods.

\begin{definition}\label{pf:open0a:def:mainx}
For a family character $f$, the seesaw transport relations $\mathrm{STR}(f)$ are the statement: for every finite set
$\mathcal I$, tuples $t_i$, constants $\alpha_i\in\C$, ramified columns $\psi_i\in\mathrm{Ram}_0(f_{t_i})$ and top forms
$G_i\in\mathrm{CT}(f_{t_i})$ ($i\in\mathcal I$),
\[
\sum_{i}\alpha_i\,c_{\mathrm{top}}^{t_i}\,\mathrm{tail}^{t_i}\bigl(\mathrm{up}^{t_i}\psi_i,G_i\bigr)=0
\ \Longrightarrow\
\sum_i\alpha_i\,c_{\mathrm{bot}}^{t_i}\,\mathrm{Per}^{t_i}(\psi_i,G_i)(y,v)=0\quad\text{for all }y,v.
\]
\formal{GlobalTheory.SeeSaw.MainX}
\end{definition}

By \cref{pf:open0a:str}, under \ref{pf:open0a:P2} and \ref{pf:open0a:P4} the family laws imply the seesaw transport
relations.

\begin{lemma}\label{pf:open0a:vw}
Let $\mathcal{K}$ be a number field, $\sigma\colon\mathcal{K}\to\C$ an embedding, $w$ a finite place of $\mathcal{K}$,
$\mathcal S$ a set of primes with $w\nmid p$ for all $p\in\mathcal S$, and $x,y\in\C$ with $\mathrm{UR}_{\mathcal S}(x,y)$.
If $a\in\mathcal{K}$ and $\sigma(a)=x/y$, then $v_w(a)=0$.
\formallabel{GlobalTheory.Fam.vw_eq_zero_of_unitRatio}\
\formallabel{GlobalTheory.T.valuation_eq_zero_of_rootOfSUnitRatio}
\end{lemma}

\begin{proof}
Let $n_0\ge1$ and $\rho\in\Q^\times$ with $(x/y)^{n_0}=\rho$, all primes dividing the numerator or the denominator of
$\rho$ lying in $\mathcal S$. Then $\sigma(a^{n_0})=\rho=\sigma(\rho)$, so $a^{n_0}=\rho$ and $a\ne0$. The numerator and
the denominator of $\rho$ are, up to sign, products of primes $p\in\mathcal S$, which are units at $w$; so they have
valuation $0$ at $w$, and $n_0\,v_w(a)=v_w(a^{n_0})=v_w(\rho)=0$. As $n_0\ge1$, $v_w(a)=0$.
\end{proof}

\begin{lemma}\label{pf:open0a:rellinks}
Let $u$ be a word, $f$ a family character for $(W,u)$ and $d$ a bottom datum. Assume that all $f_t$ are admissible,
$\mathrm{STR}(f)$, \ref{pf:open0a:L1}--\ref{pf:open0a:L6} and $\mathrm{SV}(W)$. Then there are a linear map $U$ of
$M_u$, $r'>0$ with $\|U\varphi\|=r'\|\varphi\|$ for all $\varphi$, and $c\colon\mathcal T\to\C$ with $\mathrm{UC}(c)$
and $\mathrm{Rel}_d(U,c)$.
\formallabel{GlobalTheory.MiTop.rel_of_links_alt}
\end{lemma}

\begin{proof}
Let $S,r,s_0$ be as in \ref{pf:open0a:L5}. Recall $c_{\mathrm{bot}}^t\ne0$ for all $t$ (admissibility) and
$c_{\mathrm{top}}^t\ne0$ (\ref{pf:open0a:L6}). Put
\[
U=S^{-1},\qquad r'=r^{-1},\qquad c(t)=c_{\mathrm{top}}^t\,s_0/c_{\mathrm{bot}}^t.
\]
Then $\|U\varphi\|=r^{-1}\|S(S^{-1}\varphi)\|=r^{-1}\|\varphi\|$ and $c(t)\ne0$.

\emph{Unit coefficients.} Let $w$ be a finite place of $\mathit{LL}$ and $k\in\mathit{LL}$ with
$\iota_{\mathit{LL}}(k)=c(\mu(w))/c(\mu_0)=x/y$, where $x=c_{\mathrm{top}}^{\mu(w)}/c_{\mathrm{bot}}^{\mu(w)}$ and
$y=c_{\mathrm{top}}^{\mu_0}/c_{\mathrm{bot}}^{\mu_0}$. By \ref{pf:open0a:L6} for $t=\mu(w)$,
$\mathrm{UR}_{\mathcal Q_{\mu(w)}}(x,y)$, and by the served condition $w\nmid p$ for all $p\in\mathcal Q_{\mu(w)}$.
\Cref{pf:open0a:vw} gives $v_w(k)=0$.

\emph{The relation.} Let $T\in\openzeroaRow$; by \ref{pf:open0a:L4} choose $G_0\in\mathrm{CT}(f_{\mu_0})$ with
$\mathrm{desc}^{\mu_0}(G_0)=T$, and define for $\varphi_0\in M_u$
\[
p(\varphi_0)=\frac{c_{\mathrm{top}}^{\mu_0}}{c_{\mathrm{bot}}^{\mu_0}}\,
\mathrm{tail}^{\mu_0}\bigl(\mathrm{up}^{\mu_0}(\mathrm c_0^{\mu_0}(S^{-1}\varphi_0)),G_0\bigr).
\]
Let $\varphi_0,\varphi,t,F,\alpha,\alpha'$ with $\alpha p(\varphi_0)+\alpha'c(t)\mathrm{pc}_t(\varphi,F)=0$. By
\ref{pf:open0a:L3} choose $G\in\mathrm{CT}(f_t)$ with leading component $\mathrm{lead}(t,F)$ and
$\mathrm{desc}^t(G)=\mathrm{rf}(t,F)$. By \ref{pf:open0a:L5} and \ref{pf:open0a:L2},
\[
\mathrm{tail}^t\bigl(\mathrm{up}^t(\mathrm c_0^t(S^{-1}\varphi)),G\bigr)
=s_0\,\mathrm{tail}^t\bigl(\mathrm{col}^t(\mathrm c_0^t\varphi),G\bigr)=s_0\,\mathrm{pc}_t(\varphi,F).
\]
Apply $\mathrm{STR}(f)$ to the two-element family consisting of
$(t,\ \alpha'/c_{\mathrm{bot}}^t,\ \mathrm c_0^t(S^{-1}\varphi),\ G)$ and
$(\mu_0,\ \alpha/c_{\mathrm{bot}}^{\mu_0},\ \mathrm c_0^{\mu_0}(S^{-1}\varphi_0),\ G_0)$; the columns are ramified by
\ref{pf:open0a:L1}. The top sum is
\[
\frac{\alpha'}{c_{\mathrm{bot}}^t}\,c_{\mathrm{top}}^t\,s_0\,\mathrm{pc}_t(\varphi,F)+\alpha\,p(\varphi_0)
=\alpha'c(t)\,\mathrm{pc}_t(\varphi,F)+\alpha\,p(\varphi_0)=0,
\]
so for all $y,v$,
\[
\alpha'\,\mathrm{Per}^t\bigl(\mathrm c_0^t(S^{-1}\varphi),G\bigr)(y,v)
+\alpha\,\mathrm{Per}^{\mu_0}\bigl(\mathrm c_0^{\mu_0}(S^{-1}\varphi_0),G_0\bigr)(y,v)=0.
\]
By \ref{pf:open0a:L1}, the first period is $d.\mathrm{per}(U\varphi,\mathrm{desc}^t(G))(y,v)=
d.\mathrm{per}(U\varphi,\mathrm{rf}(t,F))(y,v)$ and the second is $d.\mathrm{per}(U\varphi_0,T)(y,v)$. This is the
required relation.
\end{proof}

\begin{proposition}\label{pf:open0a:joint}
Assume \ref{pf:open0a:P2} and \ref{pf:open0a:P4}. Then $\mathrm{FC}(W)$ and $\mathrm{SV}(W)$ imply $\mathrm{JC}(W)$.
\formallabel{GlobalTheory.T.joint_of_famAllQ_of_servedQ}\ \formallabel{GlobalTheory.MiTop.joint_of_famAllQ_alt}
\end{proposition}

\begin{proof}
Take $z$ from $\mathrm{FC}(W)$. For a word $u$ let $f,d$ be as there; $d$ is pinned at $z$. The family laws contain the
admissibility of every $f_t$ and give $\mathrm{STR}(f)$ by \cref{pf:open0a:str}. \Cref{pf:open0a:rellinks} gives
$U,r',c$, and $(d,U,r',c)$ is as required.
\end{proof}


\subsubsection{The final construction}\label{sec:pf:open0a:final}
We now produce the data $(V,a,U,\zeta)$ of $\openzeroaMI$. The last two ingredients are an isometry induced by a
central element, and an anchoring of every word by a non-zero pairing; the anchoring is where the non-vanishing of
periods enters.

\begin{lemma}[Central shift]\label{pf:open0a:shift}
Assume the bottom hypotheses for $W$ and let $z\in G_f\cap Z(G)$. There is a linear map $V_0\colon H_E\to H_E$ such
that (i) $\mathrm{rf}(V_0e)(v)=R_z\bigl(\mathrm{rf}(e)(v)\bigr)$ for all $e,v$, where $R_z$ is right translation by
$z$; (ii) $\|V_0e\|=\|e\|$ for all $e$; (iii) $V_0$ preserves the true part.
\formallabel{GlobalTheory.MiTop.exists_isShift_of_center_alt}
\end{lemma}

\begin{proof}
The bottom hypotheses contain the statement that the real datum of $W.Q$ is a model: $e_E$ is injective, its image is
the good part of the multiplicity maps at the trivial level, and it is isometric when the multiplicity maps carry the
multiplicity inner product $\langle T,T'\rangle_M$ (an integral over a fundamental domain of $\Uni(Y_0)(F)$ in $G$,
with respect to a Haar measure $\nu_0$). For every level, post-composition with $R_z$ maps the good part at that level
into itself (\lref{GlobalTheory.MiTop.CentralShift.postOn_rTrL_mem_EEat_of_mem_center}: for central finite $z$,
right translation by $z$ preserves the good parts of all levels). Hence
$V_0=e_E^{-1}\circ(T\mapsto R_z\circ T)\circ e_E$ is a well-defined linear map, and (iii) holds. Since the extension of
$R_z\circ T$ to tensors is $R_z$ composed with that of $T$, (i) holds. As $z$ is central, $g\mapsto gz=zg$ is a left
translation and preserves $\nu_0$. By \lref{GlobalTheory.isMultiplicityInner_of_unit_twist_shift} (the
multiplicity inner product is unchanged when both arguments $T,T'$ are replaced by $a\,T(\cdot\,\kappa)$ and
$a\,T'(\cdot\,\kappa)$, $|a|=1$, for any $\kappa$ whose right translation preserves the Haar measure), applied with $a=1$
and $\kappa=z^{\pm1}$, $\langle R_z\circ T,R_z\circ T'\rangle_M=\langle T,T'\rangle_M$ for $T,T'$ in the good part.
Since $e_E$ is isometric, (ii) follows.
\end{proof}

A \emph{conjugate vector} is a vector $v=(v_\tau)_\tau$ of the discrete series models $M$ for which there are
conjugate-linear isometries $c_\tau$ from the model at $\tau$ onto the archimedean representation of the Hida setting
at $\tau$, intertwining the actions (through $\iota$ and $\theta$) and mapping $v_\tau$ to the distinguished vector of
the Hida setting. We write $\otimes v$ for the corresponding pure tensor in $M^{\infty}$.

\begin{lemma}[Anchor pairs]\label{pf:open0a:anchorpairs}
Assume \ref{pf:open0a:P2}, \ref{pf:open0a:P12}, the bottom hypotheses for $W$ and $\openzeroaNV(D)$. Let
$z\in G_f\cap Z(G)$ and let $v$ be a conjugate vector. Then for every word $u$:
\begin{enumerate}[label=(\alph*)]
\item every $\varphi\in M_u$ has a column form;
\item for every $T\in\openzeroaRow$ and $C\in\mathrm{Cu}(u)$ there is $b$ with $\mathrm{Pair}_u(zt_S,T(\otimes v),C,b)$;
\item there are $\varphi\in M_u$, a column form $C$ of $\varphi$, $T\in\openzeroaRow$ and $b\ne0$ with
$\mathrm{Pair}_u(zt_S,T(\otimes v),C,b)$.
\end{enumerate}
\formallabel{GlobalTheory.T.conjugate_vector_central_twist_alt}\ \formallabel{GlobalTheory.MiTop.anchor_pairs_alt}
\end{lemma}

\begin{proof}
(a) The column dictionary of $u$ maps $\mathrm{Cu}(u)$ onto $M_0(u)$ (a property of the spectral ties), and
$e_u(\varphi)\in\openzeroaram(u)\subseteq M_0(u)$.

(b) The bottom hypotheses give that $\theta$ is an archimedean family, that the level $(J,\chi)$ matches at $zt_S$, and
the gap and identity conditions of the Hida setting (\lref{GlobalTheory.T.bottom_hypothesis_matching_central_alt}:
these consequences of the bottom hypotheses hold for every central finite $z$). Under these conditions, by
\cref{pf:open0a:evalall} (which uses \ref{pf:open0a:P2} and \ref{pf:open0a:P12}),
there is a bijective linear map $j\colon\openzeroaRow\to H_X$ such that $j(T)$ is almost everywhere the complex
conjugate of $x\mapsto T(\otimes v)(\iota(x)zt_S)$, i.e.\ every cusp form of the Hida setting is the translated
evaluation of a unique element of the base row space. Then $\mathrm{Pair}_u(zt_S,T(\otimes v),C,b)$ holds with
$\phi=\overline{j(T)}$ and $b=\langle j(T),\varpi_C\rangle$.

(c) By $\openzeroaNV(D)$ choose $\varphi_r\in\openzeroaram(u)$ with $P_u(x)\ne0$, $x=\iota_u^{-1}\varphi_r$. The spectral
ties provide $C\in\mathrm{Cu}(u)$ with dictionary image $\iota_u(x)=\varphi_r$, so that $C$ is a column form of
$e_u^{-1}(\varphi_r)$, and an isometric identification of $H$ with $H_X$ under which $P_u(x)$ corresponds to
$\varpi_C$; hence $P:=\varpi_C\ne0$. With $j$ as in (b) put $T=j^{-1}(P)$; then
$b=\langle P,P\rangle\ne0$.
\end{proof}

\begin{lemma}[Anchoring a word]\label{pf:open0a:anchorword}
Let $z\in G$, $u$ a word, $d$ a bottom datum for $(W,u)$, $v$ a vector of the models, $U$ a linear map of $M_u$ with
$\|U\varphi\|=r\|\varphi\|$ for some $r>0$, and $c\colon\mathcal T\to\C$. Assume $\mathrm{BR}(z,d)$,
$\mathrm{Rel}_d(U,c)$, statement (c) of \cref{pf:open0a:anchorpairs}, and $\mathrm{rf}(t,F)\in\openzeroaRow$ for all
$t,F$. Then there are a linear isometry $R$ of $M_u$ and $p,b_0\in\C^\times$ such that for all $t$, $\varphi\in M_u$,
rows $F$, $C\in\mathrm{Cu}(u)$ and $b\in\C$:
\[
\mathrm{CF}_u\bigl(C,R(U\varphi)\bigr)\ \text{and}\ \mathrm{Pair}_u\bigl(zt_S,\mathrm{rf}(t,F)(\otimes v),C,b\bigr)
\ \Longrightarrow\ c(t)\,\mathrm{pc}_t(\varphi,F)\,b_0=p\,b.
\]
\formallabel{GlobalTheory.MiTop.AtReal.anchor_word_alt}
\end{lemma}

\begin{proof}
Let $R$ be the isometry of $\mathrm{BR}(z,d)$, and choose $\varphi_1$, a column form $C_0$ of $\varphi_1$,
$T_0\in\openzeroaRow$ and $b_0\ne0$ with $\mathrm{Pair}_u(zt_S,T_0(\otimes v),C_0,b_0)$. Since $r>0$, $U$ is injective,
hence bijective ($M_u$ is finite-dimensional); choose $\varphi_0$ with $U\varphi_0=R^{-1}\varphi_1$, so that $C_0$ is a
column form of $R(U\varphi_0)$. Let $p_0\colon M_u\to\C$ be the function given by $\mathrm{Rel}_d(U,c)$ for $T_0$, and
put $p=p_0(\varphi_0)$.

$p\ne0$: otherwise $1\cdot p_0(\varphi_0)+0=0$, and $\mathrm{Rel}_d(U,c)$ (with $\alpha=1$, $\alpha'=0$, the base tuple
and the zero row) gives $d.\mathrm{per}(U\varphi_0,T_0)(y,v')=0$ for all $y,v'$. Then $\mathrm{BR}(z,d)$, applied with
$\varphi=\varphi'=U\varphi_0$, $T=T'=T_0$, $(\alpha,\alpha')=(1,0)$, $C=C'=C_0$ and $b=b'=b_0$, gives $b_0=0$, a
contradiction.

The identity: let $t,\varphi,F,C,b$ satisfy the hypotheses, and put $\alpha=c(t)\mathrm{pc}_t(\varphi,F)$ and
$\alpha'=-p$. Then $\alpha\,p_0(\varphi_0)+\alpha'c(t)\mathrm{pc}_t(\varphi,F)=0$, so by $\mathrm{Rel}_d(U,c)$,
$\alpha\,d.\mathrm{per}(U\varphi_0,T_0)(y,v')+\alpha'\,d.\mathrm{per}(U\varphi,\mathrm{rf}(t,F))(y,v')=0$ for all $y,v'$.
Since $\mathrm{rf}(t,F)\in\openzeroaRow$, $\mathrm{BR}(z,d)$ (with $\varphi=U\varphi_0$, $\varphi'=U\varphi$, $T=T_0$,
$T'=\mathrm{rf}(t,F)$, the vector $\otimes v$, the column forms $C_0,C$ and the pairings $b_0,b$) gives
$\alpha b_0+\alpha'b=0$, i.e.\ $c(t)\,\mathrm{pc}_t(\varphi,F)\,b_0=p\,b$.
\end{proof}

\begin{lemma}[Normalisation]\label{pf:open0a:anchored}
Let $X$ be a coefficient complex with face data $n_V$, $\mathrm{Mem}$, $a$, $X_c$, $\mathrm{Pt}$, $H$, $P$, let $e_V$
and $s_{EE}$ be as in \cref{pf:open0a:def:mistar}, $V\colon E\to E$ linear, $U(u)\in M_{a(u)}(\C)$ and
$c\colon\mathrm{Mem}\times\mathcal T\to\C$. Assume:
\begin{enumerate}[label=(\roman*)]
\item $V$ satisfies \ref{pf:open0a:M1} and the $U(u)$ satisfy \ref{pf:open0a:M2};
\item for every $u$ there are $p(u),b(u)\in\C^\times$ such that for all $t,\varphi,F$
\[
c(u,t)\,\iota_0\bigl(\mathrm{Pt}(u,t)(\varphi,F)\bigr)\,b(u)
=p(u)\,\bigl\langle s_{EE}\bigl(V\Theta_t(j_R(t)(e_V(t)F))\bigr),P_u\bigl(U(u)\iota_0(\varphi)\bigr)\bigr\rangle;
\]
\item $c(u,t)\ne0$ for all $u,t$, and $\iota_0(k)=c(u,\mu(w))/c(u,\mu_0)$ implies $v_w(k)=0$, for all $u$, finite places
$w$ and $k\in\mathcal{K}$.
\end{enumerate}
Then $\openzeroaMI(X,e_V,s_{EE})$ holds, with $a(u)=p(u)/(b(u)c(u,\mu_0))$ and $\zeta(u,t)=c(u,t)/c(u,\mu_0)$.
\formallabel{GlobalTheory.MiTop.miStar2_of_anchored}
\end{lemma}

\begin{proof}
\ref{pf:open0a:M1} and \ref{pf:open0a:M2} are (i). $\zeta(u,t)\ne0$, $\zeta(u,\mu_0)=1$ and the valuation condition of
\ref{pf:open0a:M3} are (iii); $a(u)\ne0$ as $p(u),b(u),c(u,\mu_0)\ne0$. Dividing the identity (ii) by
$b(u)\,c(u,\mu_0)$ gives \eqref{pf:open0a:eq:MI}.
\end{proof}

\begin{theorem}\label{pf:open0a:final}
Assume \ref{pf:open0a:P2} and \ref{pf:open0a:P12}. Let $W$ satisfy the bottom hypotheses, assume that
$\Theta_t(j_R(t)(e_V(t)F))\in\openzeroaTrue$ for all tuples $t$ and rows $F$, $\mathrm{JC}(W)$ and $\openzeroaNV(D)$.
Then $\openzeroaMI(X_W,e_{V,W},s_{EE,W})$.
\formallabel{GlobalTheory.T.joint_star_product_nonvanishing_alt}\
\formallabel{GlobalTheory.MiTop.miStar2_at_realisation_alt}
\end{theorem}

\begin{proof}
Besides \cref{pf:open0a:shift,pf:open0a:anchorpairs,pf:open0a:anchorword,pf:open0a:anchored} the proof uses the
following statements; in \ref{pf:open0a:B2}, \ref{pf:open0a:B3} and \ref{pf:open0a:B5} the hypotheses are the bottom
hypotheses, \ref{pf:open0a:P2} and \ref{pf:open0a:P12}.
\begin{enumerate}[label=(B\arabic*),leftmargin=*]
\item\label{pf:open0a:B1} (\lref{GlobalTheory.MiTop.botHyp_match_alt}) The bottom hypotheses imply, besides the
facts used in \cref{pf:open0a:anchorpairs}, that $\mathrm{rf}(e)\in\openzeroaRow$ whenever $e\in\openzeroaTrue$.
\item\label{pf:open0a:B2} (\lref{GlobalTheory.Core.Mi.exists_conjVec}) A conjugate vector exists: the
classification \ref{pf:open0a:P12} identifies the discrete series of the Hida setting with the model $M$ by a
conjugate-linear intertwiner, using the uniqueness of discrete series with given parameter.
\item\label{pf:open0a:B3} (\cref{pf:open0a:pairform}; the \emph{pairing form}) For every
conjugate vector $v$ there is $\kappa\in\C^\times$ such that for all $e\in\openzeroaTrue$, words $u$ with member $m$,
$x\in\C^{a(m)}$, column forms $C$ of $\mathrm{col}_u(x)$ and $b$ with $\mathrm{Pair}_u(t_S,\mathrm{rf}(e)(\otimes v),C,b)$,
one has $\bigl\langle s_{EE}(e),P_m(\openzeroacj(x))\bigr\rangle=\openzeroacj(\kappa b)$. Here
$\mathrm{col}_u(x)=\sum_ix_i\,\jmath_u(\mathbf e_i)$, $\jmath_u\colon\mathit{LL}^{a(m)}\to M_u$ is the column embedding
of $u$, $\mathbf e_i$ are the unit vectors, and $\openzeroacj$ acts componentwise.
\item\label{pf:open0a:B4} (\lref{GlobalTheory.MiTop.colCrd_of_real_pack}) For every word $u$, $\mathrm{col}_u$ is
bijective, $\mathrm{col}_u(\iota_{\mathit{LL}}(x))=\jmath_u(x)$ for $x\in\mathit{LL}^{a(m)}$, and there is $r_u>0$ with
$\|\openzeroacj(x)\|_{X_c(m,v_0)}=r_u\|\mathrm{col}_u(x)\|$ for all $x$.
\item\label{pf:open0a:B5} (\cref{pf:open0a:botrel}) For $z\in G_f\cap Z(G)$ and a bottom datum $d$ pinned at $z$,
$\mathrm{BR}(z,d)$.
\item\label{pf:open0a:B6} The members of the face are the words: every $m\in\mathrm{Mem}$ is the member $m_u$ of a word
$u$; and, by the hypothesis at the package contained in the bottom hypotheses, the column pairings of $X_W$ are those
of the datum: $\iota_{\mathit{LL}}\bigl(\mathrm{Pt}(m_u,t)(\varphi,F)\bigr)=\mathrm{pc}_t(\jmath_u\varphi,F)$.
\end{enumerate}
Let $z$ and $(d_u,U_u,r_u',c_u)_u$ be given by $\mathrm{JC}(W)$, $V_0$ by \cref{pf:open0a:shift} for $z$, $v$ by
\ref{pf:open0a:B2} and $\kappa$ by \ref{pf:open0a:B3}. For every row the vector $e_{t,F}=\Theta_t(j_R(t)(e_V(t)F))$
lies in $\openzeroaTrue$ by hypothesis, and $\mathrm{rf}(t,F)=\mathrm{rf}(e_{t,F})$, so $\mathrm{rf}(t,F)\in\openzeroaRow$
by \ref{pf:open0a:B1}. For every word $u$, \ref{pf:open0a:B5} gives $\mathrm{BR}(z,d_u)$, \cref{pf:open0a:anchorpairs}
gives its statements (a)--(c), and \cref{pf:open0a:anchorword} gives $R_u$, $p_u$ and $b_{0,u}$.

Fix a word $u$ with member $m=m_u$. The map $A_u=r_u'^{-1}R_u\circ U_u$ is a linear isometry of $M_u$. Let
$U(m)\in M_{a(m)}(\C)$ be the matrix of $y\mapsto\openzeroacj\bigl(\mathrm{col}_u^{-1}(A_u\,\mathrm{col}_u(\openzeroacj
y))\bigr)$, and put
\[
c(m,t)=\openzeroacj(c_u(t)),\qquad p(m)=\openzeroacj(p_u),\qquad b(m)=\openzeroacj\bigl(\kappa\,r_u'^{-1}b_{0,u}\bigr).
\]
By \ref{pf:open0a:B4}, $\|U(m)y\|_{X_c(m,v_0)}=r_u\|A_u\,\mathrm{col}_u(\openzeroacj y)\|=r_u\|\mathrm{col}_u(\openzeroacj
y)\|=\|y\|_{X_c(m,v_0)}$, which is \ref{pf:open0a:M2}. Let $V$ be $V_0$ regarded as a map of $E$; it satisfies
\ref{pf:open0a:M1} by \cref{pf:open0a:shift}(ii),(iii). Condition (iii) of \cref{pf:open0a:anchored}: $c(m,t)\ne0$ as
$c_u(t)\ne0$; and if $\iota_0(k)=c(m,\mu(w))/c(m,\mu_0)$, then applying $\openzeroacj$ gives
$\iota_{\mathit{LL}}(k)=c_u(\mu(w))/c_u(\mu_0)$, so $v_w(k)=0$ as $c_u$ is a system of unit coefficients.

It remains to verify (ii) of \cref{pf:open0a:anchored}. Let $t$ be a tuple, $\varphi\in\mathit{LL}^{a(m)}$ and $F$ a row.
Put $\psi=\jmath_u(\varphi)\in M_u$, so that $\mathrm{col}_u(\iota_{\mathit{LL}}\varphi)=\psi$ by \ref{pf:open0a:B4}, and
$x=\mathrm{col}_u^{-1}(A_u\psi)$; then $U(m)\,\iota_0(\varphi)=U(m)\,\openzeroacj(\iota_{\mathit{LL}}\varphi)
=\openzeroacj(x)$. By \cref{pf:open0a:anchorpairs}(a),(b) choose a column form $C$ of
$\mathrm{col}_u(x)=A_u\psi=R_u(U_u(r_u'^{-1}\psi))$ and $b$ with $\mathrm{Pair}_u(zt_S,\mathrm{rf}(t,F)(\otimes v),C,b)$.
Put $e=V_0e_{t,F}\in\openzeroaTrue$. By \cref{pf:open0a:shift}(i), $\mathrm{rf}(e)(\otimes
v)=R_z\bigl(\mathrm{rf}(t,F)(\otimes v)\bigr)$, and since $z$ is central, $R_z(f)(\iota(x')t_S)=f(\iota(x')t_Sz)
=f(\iota(x')zt_S)$ for every function $f$ on $G$; hence $\mathrm{Pair}_u(t_S,\mathrm{rf}(e)(\otimes v),C,b)$, and
\ref{pf:open0a:B3} gives
\[
\bigl\langle s_{EE}(Ve_{t,F}),\,P_m\bigl(U(m)\iota_0(\varphi)\bigr)\bigr\rangle=\openzeroacj(\kappa b).
\]
\Cref{pf:open0a:anchorword}, applied to $r_u'^{-1}\psi$, gives
$c_u(t)\,\mathrm{pc}_t(r_u'^{-1}\psi,F)\,b_{0,u}=p_u\,b$; as $\mathrm{pc}_t$ is linear in the column,
$c_u(t)\,r_u'^{-1}\mathrm{pc}_t(\psi,F)\,b_{0,u}=p_ub$. By \ref{pf:open0a:B6},
$\iota_{\mathit{LL}}(\mathrm{Pt}(m,t)(\varphi,F))=\mathrm{pc}_t(\psi,F)$, and $\iota_0=\openzeroacj\circ\iota_{\mathit{LL}}$.
Therefore
\begin{align*}
c(m,t)\,\iota_0\bigl(\mathrm{Pt}(m,t)(\varphi,F)\bigr)\,b(m)
&=\openzeroacj\bigl(\kappa\,c_u(t)\,r_u'^{-1}\,\mathrm{pc}_t(\psi,F)\,b_{0,u}\bigr)
=\openzeroacj(p_u)\,\openzeroacj(\kappa b)\\
&=p(m)\,\bigl\langle s_{EE}(Ve_{t,F}),P_m\bigl(U(m)\iota_0(\varphi)\bigr)\bigr\rangle .
\end{align*}
\Cref{pf:open0a:anchored} concludes.
\end{proof}

\begin{theorem}\label{pf:open0a:mistarpunz}
Assume \ref{pf:open0a:P2}, \ref{pf:open0a:P4} and \ref{pf:open0a:P12}. If $W$ satisfies the bottom hypotheses, the level
hypothesis, $\mathrm{FC}(W)$ and $\mathrm{SV}(W)$, and $\openzeroaNV(D)$ holds, then $\openzeroaMI(X_W,e_{V,W},s_{EE,W})$.
\formallabel{GlobalTheory.T.mi_star2_of_pu_nz}\
\formallabel{GlobalTheory.MiTop.miStar2_of_bot_lev_fam_served_puNZ}
\end{theorem}

\begin{proof}
By \lref{GlobalTheory.MiTop.descent_eV_mem_level_alt}, the bottom hypotheses and the level hypothesis imply
$\Theta_t(j_R(t)(e_V(t)F))\in\openzeroaTrue$ for all $t,F$: the level hypothesis provides a unitary representation of a
finite group on $E$ compatible with the theta maps, the bijectivity at the principal level and the level types at the
two levels, and a descent argument shows that the rational rows at the true level span the intersection of the image
of $\Theta_t^*$ with the good part, so that their theta images lie in the true part. By \cref{pf:open0a:joint},
$\mathrm{JC}(W)$ holds, and \cref{pf:open0a:final} concludes.
\end{proof}

%==============================================================================
\subsection{Proofs of the deferred results}\label{sec:pf:open0a:deferred}

In this subsection we prove the ingredients which were stated and cited in the preceding subsections. The proofs
are given at the level of their main steps; the deeper ingredients are cited with a statement of what is used.

\subsubsection{Size conditions at linked bundles}
For a datum $D$ and a word $u$ let $\mathrm{Tr}_u>0$ be the \emph{level trace} of $u$ (the trace attached to the
level prime $T_u$ in the size inequalities), and for an object $O$ put $b=e^{L}$, $L=O.L.L$. The constants $C_\Gamma$,
$C_{YL}$, $C_{Tu}$ are those of the system $\kappa.C$.

\begin{lemma}\label{pf:open0a:hypL}
Statement \ref{pf:open0a:S2} holds: for all functions $A,B$ of systems of family constants, the size inequalities of $D$
at $L$ with the constants listed there hold eventually at the bundles of $(\mathcal G_p,\openzeroaLk_p)$.
\formallabel{GlobalTheory.T.degHypL_all_linked_good_members}
\end{lemma}

\begin{proof}
We combine three statements.
\begin{enumerate}[label=(\roman*)]
\item (\lref{GlobalTheory.T.trace_lower_bound_all_members'}) Eventually at the bundles of
$(\mathcal G_p,\openzeroaLk_p)$, $b^{m/2}\le\mathrm{Tr}_u$ for every word $u$.
\item (\lref{GlobalTheory.T.trace_upper_bound_all_members'}) Eventually at the bundles of
$(\mathcal G_p,\openzeroaLk_p)$, $\mathrm{Tr}_u\le b^{C_{Tu}m}$ for every word $u$.
\item (\lref{GlobalTheory.Core.HypAt.evP_hypL}) Eventually in contexts: every point $\beta$ in the context such that
$b^{m/2}\le\mathrm{Tr}_u\le b^{C_{Tu}m}$ for all words $u$ of $D_\beta$ satisfies the size inequalities of the
statement. Its proof takes $L=x^{\varepsilon}$, derives the global rows of the size data of a point
(\lref{GlobalTheory.Diag.PG.Model0.globalRows_of_sizeData}) and the node parameters
(\lref{GlobalTheory.Diag.PG.NodeParameters.nodes_of_parameters}), and deduces the size inequalities from the global
rows, the nodes and the trace bounds (\lref{GlobalTheory.T.hypotheses_of_global_rows_and_nodes}), comparing the
rank numbers with powers of $x^\varepsilon$.
\end{enumerate}
By \cref{pf:open0a:rowEcalc}(a) (whose proof applies verbatim to a conjunction of an eventual statement at the
bundles and a statement eventual in contexts), (i)--(iii) hold simultaneously eventually. At a linked bundle over a
good point $\beta$ where they hold, (i) and (ii) give the trace bounds required in (iii), and (iii) gives the size
inequalities.
\end{proof}

\begin{lemma}\label{pf:open0a:sizes}
Statement \ref{pf:open0a:S1} holds.
\formallabel{GlobalTheory.T.exists_constants_for_parameters_scale_bundle}
\end{lemma}

\begin{proof}
By \cref{pf:open0a:rowEcalc}(a) the following eight properties hold simultaneously eventually at the bundles of
$(\mathcal G_p,\openzeroaLk_p)$: the ties of $D$ (\lref{GlobalTheory.Fam.MemberSizes.ties_at_alt}); the mid-volume
formula (\lref{GlobalTheory.T.familyMidVolume_at_linked}); the size inequalities of \cref{pf:open0a:hypL} with
$A(C)=F(C_{YL}+1,C_{Tu})$ and $B(C)=G(C_{YL}+1,C_{Tu})$; the interlacing and congruence conditions
(\lref{GlobalTheory.T.dataNW_all_linked_good_members}); the trace frame
(\lref{GlobalTheory.T.type_frame_all_members'}); the rank identity at every word
(\lref{GlobalTheory.T.rankMu_all_linked_members}); the cubic node gap
(\lref{GlobalTheory.T.bigGap_of_isGoodE_isLinked}); and $H(C_{YL}+1,C_{Tu})\le L$, $C_\Gamma(O.L)=C_\Gamma$ and
$C_\Gamma\ge1$ (\lref{GlobalTheory.Fam.MemberSizes.linked_scale_ge_H_and_cgam_eq}). At a bundle where they hold
take $C_g=C_{YL}+1$, $C_T=C_{Tu}$, $C'=\max(F(C_g,C_T),8C_\Gamma)$ and $C_l=\max(G(C_g,C_T),40)$. Then
$F(C_g,C_T)\le C'$, $G(C_g,C_T)\le C_l$ and $H(C_g,C_T)\le L$; the standing size conditions hold (their five components
are among the eight properties); $O.L$ has scale parameter $L$ and $1\le C_\Gamma(O.L)\le C'/8$; the rank identity and
the cubic node gap hold. Finally the size inequalities contain $C_l\le m$, so $m\ge40$, and $r_S=m/2-1\ge1$.
\end{proof}

\subsubsection{Leading reads}

\begin{lemma}\label{pf:open0a:readat}
Let $q$ be a weighted point and $b$ a bundle over $(D_q,O_q,c)$ with a bundle witness, and assume that the reading data
of $q$ admit a reading structure. Then for every word $u$ there is a top map $j$ such that the leading read of
$\mathrm{pk}(b)$ holds for $(u,j)$ at the identity: $j$ is admissible for $u$, and for every tuple $t$ the level
condition holds at the level data of $t$ and the coefficients of every row of $t$ are read by $j$.
\formallabel{GlobalTheory.T.exists_read_at_packB_of_bWit}
\end{lemma}

\begin{proof}
From the bundle witness we obtain the admissibility of the construction, pins, standard columns and the standard twist
condition (\lref{GlobalTheory.Fam.colsStd_of_bWit_alt}); with these, the forms at the level of the type places
have Gram matrix $1$ in frame coordinates (\lref{GlobalTheory.T.trueColumn_spaces_standard_in_frameCoords}).
The standardness of the core contained in the reading structure gives a semilinear \emph{valuation map} from the
values of the base to the coefficient space of the package, which satisfies the valuation equivalence and the
package valuation property (\lref{GlobalTheory.T.exists_valEqv_of_bWit_std}), and the witness gives a vanishing
property of the package at the identity (\lref{GlobalTheory.T.pk_at_vanishing_of_is_uni}). From these,
\lref{GlobalTheory.TopVal.exists_is_j_pair_val_of_vanishing} gives a top map $j$ which is admissible for $u$,
together with the values of $j$ on the pairs of each tuple. Now let $t$ be a tuple. The level condition at the
identity holds: the level data of $t$ are those of the reference tuple
(\lref{GlobalTheory.T.level_group_reference_tuple_alt}), an element lies in the level group if and only if it
lies in the twisted class (\lref{GlobalTheory.T.mem_level_iff_twisted_class_alt}), and the level character is $1$
on the twisted class (\lref{GlobalTheory.T.level_character_twisted_class_alt}); so $1$ lies in the level group
with character value $1$. Finally, for every row $F$ of $t$ the reading structure provides a one-form $\varphi$ with the
reading identity at $u$ (\lref{GlobalTheory.T.one_form_exists_heq_alt}), and
\lref{GlobalTheory.TopVal.reads_pk_alt} turns this identity and the values of $j$ into the reading of the
coefficients of $F$ by $j$.
\end{proof}

\subsubsection{Surjectivity along chains}

\begin{proposition}\label{pf:open0a:ontoclassical}
Assume \ref{pf:open0a:P1}--\ref{pf:open0a:P5}. If a datum $D$ satisfies the Rallis inner product formula of
\cref{pf:open0a:onto}(a) and the parameter transfer of \cref{pf:open0a:onto}(b), then $\mathrm{Onto}(D)$.
\formallabel{GlobalTheory.T.thetaLiftChain_onto_cuspForms_of_innerProduct}
\end{proposition}

\begin{proof}
The proof establishes a sequence of properties of the chains of $D$, each from the previous ones.
\begin{enumerate}[label=(\arabic*),leftmargin=*]
\item The level facts (\lref{GlobalTheory.T.thetaChain_level_properties}) and, by \ref{pf:open0a:P1}, the
finite dimensionality of the cusp spaces of all localisations
(\lref{GlobalTheory.T.thetaChainCuspSpaceFiniteDim_of_harishChandra}).
\item By \ref{pf:open0a:P2}, the Petersson inner products on these spaces exist
(\lref{GlobalTheory.T.thetaChainInnerProductsExist_of_rapidDecay}); the multiplicity spaces are spanned by Hecke
eigen-elements (\lref{GlobalTheory.T.multiplicity_space_spanned_by_satake_eigen_elements}); the Hecke parameters
of distinct eigen-elements are linearly independent
(\lref{GlobalTheory.T.thetaChain_heckeParameters_linearIndependent}); the Hecke parameter of an element is
determined by one non-zero value (\lref{GlobalTheory.T.hecke_parameters_of_element_from_one_nonzero_value}); and
the occurring parameters satisfy the half-rank norm bound
(\lref{GlobalTheory.T.theta_chain_occurring_params_norm_half_rank_bound}).
\item The Fock vectors of the chain satisfy the level laws (\lref{GlobalTheory.T.thetaLiftChain_fockVector_levelLaws}),
and their coefficients vanish off the level (\lref{GlobalTheory.T.thetaChainCoef_vanishes_off_level}).
\item By \ref{pf:open0a:P3} and \ref{pf:open0a:P4}, the archimedean theta correspondence holds on the towers of real
data (\lref{GlobalTheory.T.archimedeanThetaCorrespondence_towers_real_data}); hence the archimedean zeta integrals
do not vanish (\lref{GlobalTheory.T.thetaChain_realZeta_nonvanishing}), and with (2) and (3) the local zeta
integrals of the chain do not vanish (\lref{GlobalTheory.T.thetaChainLocalZetaNonvanishing_of_real}).
\item By the inner product formula and (4), the upward lifts of non-zero eigen-elements are non-zero
(\lref{GlobalTheory.T.thetaLiftUpwards_ne_zero_of_lValue}): the Petersson norm of the lift is a non-zero multiple of
a non-vanishing $L$-value times the local zeta integrals.
\item By \ref{pf:open0a:P2}--\ref{pf:open0a:P5}, the upward lifts lie in the cusp spaces
(\lref{GlobalTheory.T.thetaLiftUpwards_in_cuspSpan}) and have adjoints for the Petersson inner products
(\lref{GlobalTheory.T.thetaLiftUpwardsAdjoint_of_rapidDecay}).
\item From (1), (2), (5), (6) and the parameter transfer: at each step the adjoint of the upward lift does not vanish on
the image of an eigen-element (\lref{GlobalTheory.T.thetaLiftChain_adjoint_nonvanishing_step}); with the
non-vanishing of the $L$-values of the chain (\lref{GlobalTheory.T.thetaChain_LValues_nonvanishing}) and the
parameter bound of (2), the same holds for the composite chains (\lref{GlobalTheory.ChainOnto.chainUp_of}).
\item The cusp space at level $0$ is spanned by eigen-elements with linearly independent parameters, and by (7) the
composite of the reachable theta maps does not vanish on any of them; hence every reachable theta map is surjective
(\lref{GlobalTheory.T.thetaLiftChain_surjective_of_adjoint_nonvanishing}). This is $\mathrm{Onto}(D)$.\qedhere
\end{enumerate}
\end{proof}

\begin{lemma}\label{pf:open0a:linkonto}
Assume the package hypotheses for $W$ and $\mathrm{Onto}(D)$. If a family character $f$ for $(W,u)$ satisfies its row
conditions, then the surjectivity \ref{pf:open0a:L4} holds.
\formallabel{GlobalTheory.T.linkOnto_of_ontoAll'}
\end{lemma}

\begin{proof}
The package defines a \emph{base chain} of $D$ with localisations $Y_k$ ($k\le m$, the stages of the spaces $Y_k$) and a
localisation $X'$ at level $m$; the package hypotheses give that the base chain and these localisations are admissible
(\lref{GlobalTheory.T.baseChain_com_z_is_of_pkHyp}, \lref{GlobalTheory.T.baseChain_stageY_is_of_pkHyp_of_lt},
\lref{GlobalTheory.T.baseChain_stageX_is_of_pkHyp}). Let $T\in\openzeroaRow$. Under the row conditions, the
descents of the base tuple chain $f_{\mu_0}$ from level $k$ are theta maps reachable in $k$ steps from $Y_0$
(\lref{GlobalTheory.BaseChain.reach_stageY_dn}), and the transition from $Y_{m-1}$ to $X'$ is an admissible step
(\lref{GlobalTheory.BaseChain.chainX_is_succ_succ}); so the descent $\mathrm{desc}^{\mu_0}$ is a theta map
reachable in $m$ steps. The base row space lies in the cusp space of $Y_0$
(\lref{GlobalTheory.T.mem_stageY_zero_cc_of_mem_zero_set}), and the package hypotheses contain the node gap
condition. By $\mathrm{Onto}(D)$ the descent is surjective, so $T=\mathrm{desc}^{\mu_0}(G_0)$ for some top form $G_0$.
\end{proof}

\subsubsection{Geometric families and placement}

\begin{proposition}\label{pf:open0a:geoall}
Assume \ref{pf:open0a:P3}, \ref{pf:open0a:P4} and \ref{pf:open0a:P13}, and $\mathcal A(W)$ for a real package over an
extra datum of $D=I.\mathrm{rec}(l,Z)$. Then for every word $u$ there is a geometric family $a$ for $(W,u)$ which is
admissible and true, whose transfer matrix at level $m$ is the identity, and whose bottom Bessel datum $a.d.B$ is
framed (statement \ref{pf:open0a:F2}).
\formallabel{GlobalTheory.T.all_true_family_datum_alt}
\end{proposition}

\begin{proof}
By strong approximation \ref{pf:open0a:P13} (\lref{GlobalTheory.T.all_true_family_geometry_alt}) there are a
\emph{geometric chain} (points $x_j$ with lattices at the levels $j\le m$) and constants $c_0(j)$ such that: the points are
admissible; the points are pinned (the lattices of $x_j$ are the standard integral lattices); the top element
$b_T$ is finite; the transfer matrix at level $m$ is the identity; and the line, embedding, orthogonality and local
conditions of geometric chains hold. The geometric conditions of $W$ give a leading read at level one for $u$. Then
\lref{GlobalTheory.GeoExists.core_data_of_point} (which uses \ref{pf:open0a:P3}, \ref{pf:open0a:P4}, the package
hypotheses and the level condition of $W$, the non-vanishing of $T_u$ and the fact that neither the type places nor the
primes of the sets $Q_t$ lie over $T_u$) completes this chain, the leading read and the column data at the points
(\lref{GlobalTheory.T.all_true_family_geometry_column_alt}) to a geometric family $a$ which is admissible, whose
geometric chain is the given one, whose bottom Bessel datum is framed, and which satisfies the remaining condition of
truth. Truth of $a$ is the conjunction of the pin of the points, this condition and the embedding and $z$-conditions of
the chain; the transfer matrix is the identity by construction.
\end{proof}

\begin{lemma}\label{pf:open0a:isfam}
Assume \ref{pf:open0a:P3}, \ref{pf:open0a:P4} and the package hypotheses for $W$. If a family character $f$ for
$(W,u)$ satisfies its row conditions, then its commuting chain $f.c$ is admissible; if moreover $f$ satisfies its
column conditions, then every tuple chain $f_t$ is admissible (statement \ref{pf:open0a:F6}).
\formallabel{GlobalTheory.T.famCh_is_and_tuples_is_of_colOK}
\end{lemma}

\begin{proof}
The admissibility of $f.c$ consists of: the standardness of the additive character and the archimedean family (part of
the bottom hypotheses); the conditions on the points contained in the row conditions; the co-compactness of
$\Uni(X')(F)\backslash\Uni(X')(\A)$, which holds as $X'$ is totally positive definite
(\lref{GlobalTheory.Unf.CoCptOf.coCpt_of_isAniso_or_posDef}); the measure conditions contained in the row
conditions; and the spanning statements: at each level $j+1<m$ the images of the polynomial intertwiners from the model
at level $j+1$ span the smooth tensors of the model at level $j$ (\lref{GlobalTheory.T.itSpans_MY_succ}), and
similarly from $X'$ to the last level (\lref{GlobalTheory.T.itSpans_MY_last_MX}); both use \ref{pf:open0a:P3}
and \ref{pf:open0a:P4}. For the tuple chains: the row choices provide the forward and backward step data; the step
vectors are theta functions, so their images are pure tensors of the model; the lifts of cusp forms from level $j+1$
equal $c_I(j)$ times the theta maps of the chain at every level, including the last
(\lref{GlobalTheory.T.lift_s_eq_ci_smul_io}, \lref{GlobalTheory.T.lift_s_last_eq_ci_succ_smul_io}); the
constants $c_I(j)$ are non-zero by the measure conditions; and the kernel, $e$-, stabiliser and support conditions are
the column conditions.
\end{proof}

\begin{proposition}\label{pf:open0a:placed}
Assume $\mathcal A(W)$. Let $a$ be an admissible and true geometric family for $(W,u)$ and let $\mathrm{row},\mathrm{col}$
be choices for which the commuting chain of $f=a.\mathrm{fam}(\mathrm{row},\mathrm{col})$ is admissible. Then the placement
conditions hold for $(f,a.d.B)$ (statement \ref{pf:open0a:F7}).
\formallabel{GlobalTheory.T.placed_p_of_true_geometry}
\end{proposition}

\begin{proof}
The placement conditions consist of the local field and type datum, the commuting vector datum, the local group data and
their indexing by the type places, all taken from $a$, together with: the admissibility of $f.c$ (hypothesis) and of the
vector datum, its tie with $f$, the shape condition and the pin of the local groups (from the admissibility of $a$); the
local group reading (from the tie); the \emph{primitivity at the top}; and the \emph{local bottom compatibility}.

Primitivity at the top: it is required at the primes $p$ of $F$ outside a set which, by the tie, consists of the primes
containing $T_u$; so let $p$ be a prime with $T_u\notin p$. The point of the true chain at the top level is pinned, and the distinguished vector there has
coordinates $(0,\dots,0,\pi_u)$, where $\pi_u\in\OO_E$ is the element of the construction attached to $u$, with
$\pi_u\bar\pi_u=T_u$; in particular all its coordinates are integral. A vector with integral coordinates one of which,
$x_i$, satisfies $x_i\,c=T_u$ for some integral $c$, where $T_u\notin p$, is primitive at $p$ (a common divisor of the
coordinates divides $T_u$). Applying this with $x_i=\pi_u$ and $c=\bar\pi_u$ gives the primitivity.

Local bottom compatibility: the level of $a.d.B$ and its local data are the bottom data of $a$ (bottom clauses of the
admissibility); the norm condition holds with scalar $1$ and the constant $-\xi_0^{m/2}$ (bottom clause and norm
condition of $a$); the local groups are standard; the lattice of the chain equals the base lattice off the twist
(\lref{GlobalTheory.Fam.lattice_eq_base_off_twist_alt}); and for every element $\kappa$ of the finite adelic
group at level $0$ the local data and the reduction level at $\kappa$ are those required
(\lref{GlobalTheory.T.true_geometry_local_data_alt}, \lref{GlobalTheory.T.true_geometry_reduction_level_alt}).
\end{proof}

\subsubsection{Column choices}

\begin{proposition}\label{pf:open0a:colside}
Assume \ref{pf:open0a:P3}, \ref{pf:open0a:P4} and $\mathcal A(W)$. Let $a$ be an admissible and true geometric family
for $(W,u)$ whose transfer matrix at level $m$ is the identity, $\mathrm{row}$ admissible row choices which are pairwise
of the same shape, and $k$ a system of family constants for $a$. Then there are column choices $\mathrm{col}$, pinned
tuples $w_0(t)$, rescaled tuples $w(t)$ and families $F_t$ ($t\in\mathcal T$) such that the column condition holds for
$a.\mathrm{fam}(\mathrm{row},\mathrm{col})$ and the column output condition holds for $(k,w_0,w,F)$ (statement
\ref{pf:open0a:F4}).
\formallabel{GlobalTheory.T.true_geometry_family_invariance_alt}
\end{proposition}

Here the \emph{column condition} asks for constants $c_0$, the pin of the points of $a$, and for every tuple $t$ a
pinned chart of the columns of $t$ (characters at each level with the twist datum and its reading) and a top element
with the embedding, level, top and total-constant conditions. The \emph{column output condition} asks, for every tuple
$t$: the family laws of $F_t$ at the pinned tuple $w_0(t)$, the level identity, the vector identity, the value and
$e$-conditions at each level, the unit ratio of the total constants, and the vanishing of the zero-level terms.

\begin{proof}
For every column choice the family $a.\mathrm{fam}(\mathrm{row},\mathrm{col})$ satisfies its row conditions (from the
admissibility of $a$ and the hypotheses on $\mathrm{row}$), so its commuting chain is admissible by \cref{pf:open0a:isfam}.

\emph{Data.} The following data are provided by results of the development, for the given $a$:
\begin{enumerate}[label=(\alph*)]
\item for every tuple $t$, a \emph{quadratic extension datum} $J_t$ which reads the family
(\lref{GlobalTheory.T.exists_quadratic_extension_reads_of_true}); each $J_t$ has a unique character at level $m$
(the product of the cardinalities of its character sets at level $m$ is $1$);
\item constants $C_j$ and, for every $t$, scalars $s_t$, roots of unity $\zeta_t$ and constants $\rho_t$ satisfying the
\emph{row-pair relation} of the levels (\lref{GlobalTheory.T.true_geometry_row_pair_relation_alt_alt_alt}; this uses
that the transfer matrix is the identity);
\item non-zero volume constants $V_j$ ($j<m$) and $V_T$ with the volume identities at each level
(\lref{GlobalTheory.T.true_geometry_exists_volume_data_alt});
\item a kernel datum at each level (\lref{GlobalTheory.T.true_geometry_exists_kernel_alt});
\item finite combinations $d_{t,j,\alpha,\gamma}$ realising the twisted local values of $J_t$ at the levels $j<m$
(\lref{GlobalTheory.T.true_geometry_twist_local_value_alt});
\item a joint twist relation $S'_t$ (\lref{GlobalTheory.T.true_geometry_joint_twist_relation_alt}).
\end{enumerate}

\emph{Construction.} Put, for $t\in\mathcal T$ and $j<m$,
\begin{gather*}
c_Q(t,j)=\prod_i\bigl(\mathrm{idx}_i(j+1)\,\mathrm{idx}_i(j+2)\bigr)^{-1},\qquad
c_K(t,j)=V_j\,C_j\,c_0(j)\,\rho_t(j),\\
c_P(t)=V_T\sum_\beta s_t(m,\beta),\qquad
c_V(t,j)=\kappa_j\,k.\mathrm{sc}(t,j)\,k.\mathrm{sc}(t,j+1)\,c_Q(t,j)\,c_G(j)\,c_T(t,j),
\end{gather*}
where $\mathrm{idx}_i$ are the indices of the twist data of $J_t$, $c_T(t,j)$ is the monomial value of the type datum of $a$
at the constants $(k.r_j,k.e_j)$ and the character $\mu_t(j+1)$, $c_G$ is the constant of $k$, and $\kappa_j$ is the
constant of $k$ attached to $j$ when $j+2(\sigma+1)+2=m$ for some $\sigma$ (and $\kappa_j=1$ otherwise). Let $\mathrm{col}$ be
the column choice determined by $a$, the data $J$, the kernel datum, the combinations $d$, the relations $S'_t$ evaluated
at the constants $k.\varepsilon_j$, the scalars $s$ and the constants $c_V,c_K,c_P$; let $F_t$ be the family of the row
and column choices at $t$, $w_0(t)$ its pinned tuple with the constants $k.\mathrm{sc}(t,\cdot)$ and $c_T(t,\cdot)$, and
$w(t)$ the rescaling of $w_0(t)$ by $s_t$.

\emph{Verification of the column condition.} The pinned chart of $t$ has weights $C_j\,\mathrm{wt}_t(j,\alpha)\,\zeta_t(j,
\alpha,\gamma)$: its admissibility and liveness come from (a), its characters from (a) and (d), its rational data from (e)
and (b), its constants from (c) and (b) (the identity $V_jC_jc_0(j)\rho_t(j)=\dots$ is the volume identity combined with
the row-pair relation), and its ramification and reading conditions hold by construction. The top element uses the unique
character $\beta_0$ of (a) and the level element provided by (c), with total constant $s_t(m,\beta_0)\,\mathrm{wt}_t(m,\beta_0)$.

\emph{Verification of the output condition.} The prime sets of $w_0(t)$ are the primes of $F$ below the primes of $Q_t$
(\lref{GlobalTheory.T.true_geometry_twist_local_data_alt}), and the pins of $F_t$ hold by construction. The column
conditions of $a.\mathrm{fam}(\mathrm{row},\mathrm{col})$ hold: the vacuum conditions at consecutive levels follow from the
pins of the points, the kernel datum and (e); the data conditions from the joint twist data
(\lref{GlobalTheory.T.true_geometry_joint_twist_data_alt}); the stabiliser conditions from the fact that the chain
kernels preserve the ramified forms of the live charts
(\lref{GlobalTheory.T.chainKernel_preserves_ramifiedForms_liveCharts}). The $\kappa$-equations of $k$ and $w(t)$ hold
by the choice of $c_V$. If $x,y$ are roots of unity and $c\ne0$, then $(cx)/(cy)$ is a root of unity, so
$\mathrm{UR}_{\mathcal S}(cx,cy)$ holds for every $\mathcal S$; this gives the unit ratio of the total constants, since
$\sum_\beta s_t(m,\beta)=s_t(m,\beta_0)$. The family laws of $F_t$ are \lref{GlobalTheory.T.true_geometry_family_final_alt}
(about $1.7\cdot10^4$ lines; it derives the seesaw laws at the twist places from a rescaling by roots of unity), the level
identity is \lref{GlobalTheory.T.true_geometry_quadratic_twist_level_alt}, the vector identity is
\lref{GlobalTheory.T.true_geometry_joint_twist_identity_alt}, the value and $e$-conditions are
\lref{GlobalTheory.T.true_geometry_joint_twist_value_alt} and
\lref{GlobalTheory.T.true_geometry_joint_twist_data_alt}, and the vanishing of the zero-level terms is
\lref{GlobalTheory.T.true_geometry_joint_twist_bound_alt}.
\end{proof}

\begin{proposition}\label{pf:open0a:lawslinks}
Assume \ref{pf:open0a:P3}, \ref{pf:open0a:P4}, \ref{pf:open0a:P13} and $\mathcal A(W)$. Let $a$ be an admissible and true
geometric family for $(W,u)$ and $\mathrm{row}$ admissible row choices which are pairwise of the same shape, and assume
that the conclusion of \cref{pf:open0a:colside} holds for every system $k$ of family constants for $a$. Then there are
column choices $\mathrm{col}$ such that $f=a.\mathrm{fam}(\mathrm{row},\mathrm{col})$ satisfies its row conditions,
$\mathrm{Laws}(f)$, the chart condition at level $0$, \ref{pf:open0a:L2}, \ref{pf:open0a:L5} and \ref{pf:open0a:L6}
(statement \ref{pf:open0a:F5}).
\formallabel{GlobalTheory.T.true_geometry_family_column_data_alt}
\end{proposition}

\begin{proof}
\emph{The constants.} Apply the hypothesis to the trivial system of constants to obtain some column choice $\mathrm{col}_0$;
by \cref{pf:open0a:isfam} the commuting chain $c$ of $a.\mathrm{fam}(\mathrm{row},\mathrm{col}_0)$ (which does not depend
on the column choice) is admissible. The finite integral laws of the seesaw for the lattice frame pins
(\lref{GlobalTheory.T.seeSaw_latticeFramePins_finiteIntegralLaws}; about $2.1\cdot10^4$ lines) give units
$\varepsilon_j$, bilinear operators $\mathrm{op}_j$, exponents $r_j,e_j$ and constants $c_G(j)$ with the corresponding
integral laws, $\kappa$-relations and admissibility of the vector datum. The non-vanishing of scalars up to roots of unity
(\lref{GlobalTheory.T.nonzero_scalars_roots_of_unity_alt}) gives non-zero scalars $s_j$ and roots of unity
$\zeta_t(j)$ with the vector and twist properties. Let $k$ be the system of family constants
$(\varepsilon,\mathrm{op},r,e,c_G,\ \mathrm{sc}(t,j)=s_j\zeta_t(j))$, and let $(\mathrm{col},w_0,w,F)$ be given by the
hypothesis for $k$; put $f=a.\mathrm{fam}(\mathrm{row},\mathrm{col})$.

\emph{Consequences of the output condition.} The prime sets of $w(t)$ are the primes below the primes of $Q_t$. The vector
function condition of $f$ at $w(t)$ with the constants $k.\mathrm{sc}(t,\cdot)$ follows from the property of $\zeta$; the
tuple chains are admissible by \cref{pf:open0a:isfam}; the Weil conditions and the characters of the local groups hold
(\lref{GlobalTheory.T.weilOK_of_rank_one_gram'}, using that the bottom lattice has rank one with the prescribed
Gram matrix); and the quadratic family level identity holds (\lref{GlobalTheory.T.quadratic_family_level_identity_alt}).
These give the vector forms of $k$ at every $w(t)$.

\emph{The laws.} The class pin data of the tuple chains of $f$ hold (\lref{GlobalTheory.T.cls_pin_d_of_true_family};
about $8.8\cdot10^3$ lines). $\mathrm{Laws}(f)$ is then the conjunction of: admissibility of $c$ and of the $f_t$; local
finiteness; the reality equations (\lref{GlobalTheory.T.realEq_of_colOK'}); the valuation and last-step laws (from
the vector forms and the finite integral laws); the kernel law, obtained from the definite class data
(\lref{GlobalTheory.T.seeSawKernelLiftCombination_of_definiteClassData}); the period law, the factorisation of the last
periods (\lref{GlobalTheory.T.seeSawLastPeriodFactorization_of_classData}); and the reading law, from the continuity
and automorphy of the scalar forms.

\emph{The links.} The chart condition at level $0$ is the zero-level part of the output condition. The tail identity
\ref{pf:open0a:L2} follows from the leading read of $a$ (\lref{GlobalTheory.MiTop.linkTail_of_topReadAt_alt}), since
at the base tuple the level-$0$ embedding of a column is its identification with the ramified subspace. The commutation
\ref{pf:open0a:L5} is the unit-shift relation of the type data
(\lref{GlobalTheory.T.type_data_unit_shift_relation_alt}), whose hypotheses are: the ramified units at the type
places, the kernel back-transport of the tuple chains
(\lref{GlobalTheory.T.seeSawTuple_kernelBack_of_classKernelData}), the tail reading of the last period classes
(\lref{GlobalTheory.T.lastPeriodClasses_give_tailReading}), and the commutation of the definite lift with the
level-prime operators (\lref{GlobalTheory.T.definiteLift_comm_levelPrimeOperators}). For the ratio condition
\ref{pf:open0a:L6}: the scalars of $w(t)$ and $w(\mu_0)$ differ by the root of unity $\zeta_t(j)/\zeta_{\mu_0}(j)$, so their
quotients are roots of unity and in particular $\mathcal Q_t$-unit ratios; the constants $c_T$ are unit ratios up to roots
of unity (\lref{GlobalTheory.T.root_of_unity_theta_twist_alt}); the constants $c_Q$ are compared with the products
$c_I(j)c_I(j+1)$, which are non-zero (\lref{GlobalTheory.T.cI_mul_prod_idx_eq_cI_mu0}); and the total constants
$c_{\mathrm{top}}^t/c_{\mathrm{bot}}^t$ are unit ratios by the factor-ratio lemma
(\lref{GlobalTheory.T.seeSaw_totalConstant_ratio_of_factor_ratios}), which also gives $c_{\mathrm{top}}^t\ne0$.
\end{proof}

\subsubsection{The central pin}

\begin{proposition}\label{pf:open0a:botdat}
If $W$ satisfies the package hypotheses and $r_S\ge1$, there is $z\in G_f\cap Z(G)$ such that for every word $u$, every
family character $f$ for $(W,u)$ and every framed Bessel datum $B$ satisfying the placement conditions for $(f,B)$, there
is a bottom datum $d$ pinned at $z$ with $d.B=B$ and $d.b=f.g.b(0)$ (statement \ref{pf:open0a:F1}).
\formallabel{GlobalTheory.T.exists_botDat_of_placedP_framed}
\end{proposition}

\begin{proof}
For a family character $f$, a Bessel datum $B$ and $a_0,a\in\Uni(X_0)(\A)$ put
$\tau(f,B,a_0,a)=\iota(1,x_u)^{-1}\,\mathrm{emb}_B(a_0a)\,f.g.b(0)\in G$, the \emph{translation element}. We say that $B$
is a translate at $(a_0,a)$ if $B$ is adapted, $a_0$ is an adelic transport correction, $a$ is finite and central, and the
level $(J,\chi)$ matches at $\tau(f,B,a_0,a)$. If there is no triple $(u,f,B)$ with placement conditions and $B$ framed,
take $z=1$. Otherwise fix such a triple $(u_0,f_0,B_0)$.

(i) By \lref{GlobalTheory.T.exists_isAt_of_placed_framed} (about $7.4\cdot10^3$ lines), every placed framed $B$ is a
translate at some $(a_0,a)$; let $B_0$ be a translate at $(a_0,a)$.

(ii) The bottom hypotheses give that the level matches at $t_S$. By
\lref{GlobalTheory.T.matching_translates_bottom_alt} (about $3.5\cdot10^3$ lines; it uses $r_S\ge1$), there are
$z\in G_f\cap Z(G)$, a finite central $c\in\Uni(X_0)(\A)$ and $k_0\in J$ with
$\iota(1,x_{u_0})^{-1}\,\mathrm{emb}_{B_0}(c)^{-1}\,\iota(1,x_{u_0})\cdot\tau(f_0,B_0,a_0,a)=z\,t_S\,k_0$. Since $c$ is
central, $a_0(c^{-1}a)=c^{-1}(a_0a)$, and since $\mathrm{emb}_{B_0}$ is a homomorphism, a direct computation in $G$ gives
$\tau(f_0,B_0,a_0,c^{-1}a)=z\,t_S\,k_0$; as $c^{-1}a$ is finite and central and the level matches at $z\,t_S\,k_0$
(matching at $zt_S$ is preserved by right multiplication by elements of $J$), $B_0$ is a translate at $(a_0,c^{-1}a)$.

(iii) Let $(u,f,B)$ be placed and framed. By (i), $B$ is a translate at some $(b_0,b)$. The compatibility of placed Bessel
data (\lref{GlobalTheory.T.placed_bessel_compatibility_alt}; about $4.5\cdot10^3$ lines; it uses $r_S\ge1$) gives
$a_1$ such that $B$ is a translate at $(b_0,a_1)$, and $k\in J$ with
$\tau(f,B,b_0,a_1)=\tau(f_0,B_0,a_0,c^{-1}a)\,k=z\,t_S\,k_0k$. Put $d=(B,f.g.b(0),b_0,a_1,k_0k)$. Then $d.B=B$,
$d.b=f.g.b(0)$, and $d$ is pinned at $z$: $B$ is adapted (placement conditions), $b_0$ is a transport correction and $a_1$
is finite and central (translate), $k_0k\in J$ ($J$ is a group), and the identity
$\mathrm{emb}_B(b_0a_1)\,f.g.b(0)=\iota(1,x_u)\,z\,t_S\,k_0k$ is the equation $\tau(f,B,b_0,a_1)=zt_Sk_0k$.
\end{proof}

\subsubsection{The seesaw transport relations}

\begin{proposition}\label{pf:open0a:str}
Under \ref{pf:open0a:P2} and \ref{pf:open0a:P4}, $\mathrm{Laws}(f)$ implies $\mathrm{STR}(f)$ for every family character
$f$ (more generally, for every family of tuple chains of an admissible commuting chain).
\formallabel{GlobalTheory.T.seesawTransportRelations_of_familyLaws}
\end{proposition}

\begin{proof}
The proof passes through three intermediate relations.
\begin{enumerate}[label=(\alph*)]
\item \emph{Double lifts.} From the admissibility of the chains, the reality equations, the kernel law and the reading law,
the double-lift combinations of the tuple chains satisfy their linear relations
(\lref{GlobalTheory.T.seeSawDoubleLiftCombination_of_laws}); the ingredients are the commutation of the two Weil
actions of a seesaw, the boundedness of the degree of central eigen-theta Fock vectors, and the summation of lifts.
\item \emph{Double steps.} By \ref{pf:open0a:P2}, values of rapidly decreasing functions are controlled; by
\ref{pf:open0a:P4}, the double-step unfolding of Bessel periods (which follows from the descent of Weil data), the
movement of values and the summation of periods, linear relations between the periods at level $j+2$ imply linear
relations at level $j$ (\lref{GlobalTheory.T.seeSaw_doubleStep_linearRelations_periods}); this uses the local
finiteness, the valuation law, (a) and the reading law.
\item \emph{The last two steps.} By \ref{pf:open0a:P2} and \ref{pf:open0a:P4}, relations between the top pairings transport
to relations between the periods two levels below (\lref{GlobalTheory.T.seeSaw_lastTwoSteps_transport_linear_relations});
this uses the last-step law, (a), the period law and the reading law.
\end{enumerate}
Finally, the transport relations from level $2(n+1)$ to level $0$ follow by composing (c) with the double steps (b)
(\lref{GlobalTheory.SeeSaw.mainX_of_stepX_lastX}).
\end{proof}

\subsubsection{Translated evaluations, the pairing form and the bottom relation}

\begin{lemma}\label{pf:open0a:evalall}
Assume \ref{pf:open0a:P2} and \ref{pf:open0a:P12}. Let the discrete series models $M$ have the parameters of level $0$, let
$\theta$ be an archimedean family, let the Hida setting satisfy the gap condition with constant $100\,r^4$ and the identity
condition, let the level $(J,\chi)$ match at $g\in G$, and let $v$ be a conjugate vector. Then $\|v_\tau\|=1$ for every
$\tau$, and there is a bijective conjugate-linear map $j\colon\openzeroaRow\to H_X$ such that $j(T)$ is almost everywhere
the complex conjugate of $x\mapsto T(\otimes v)(\iota(x)g)$.
\formallabel{GlobalTheory.T.evalAll_and_norm_eq_one_of_conjVec}
\end{lemma}

\begin{proof}
The conjugate-linear isometry $c_\tau$ maps $v_\tau$ to the distinguished vector of the Hida setting, a unit vector of the
lowest $K$-type; hence $\|v_\tau\|=1$ and $v_\tau\ne0$. For $T\in\openzeroaRow$ the translated evaluation
$x\mapsto T(\otimes v)(\iota(x)g)$ is a square-integrable automorphic function
(\lref{GlobalTheory.T.translateToSecondModel_cuspFormMap_value_automorphic}; by \ref{pf:open0a:P2}); let $f(T)\in
L^2_X$ be its class. Its complex conjugate lies in the space $H_X$ of cusp forms
(\lref{GlobalTheory.T.conj_translated_evaluation_mem_all_rows_space}; by \ref{pf:open0a:P2}). The map
$j(T)=\overline{f(T)}$ is additive and $j(aT)=\bar a\,j(T)$. It is injective: if $j(T)=0$, the translated evaluation of
$T(\otimes v)$ vanishes almost everywhere, and since all $v_\tau\ne0$ this forces $T=0$
(\lref{GlobalTheory.T.cuspFormMap_eq_zero_of_translatedEval_ae_zero}). It is surjective: every $\phi\in L^2_X$ with
$\bar\phi\in H_X$ is almost everywhere the translated evaluation of $T(\otimes v)$ for some $T\in\openzeroaRow$
(\lref{GlobalTheory.T.conj_row_eq_translated_evaluation_cusp_form_map}; by \ref{pf:open0a:P12}, through the
identification of the discrete series of the Hida setting with $M$ given by the conjugate vector).
\end{proof}

\begin{proposition}\label{pf:open0a:pairform}
Assume \ref{pf:open0a:P2}, \ref{pf:open0a:P12} and the bottom hypotheses for $W$, and let $v$ be a conjugate vector. Then
the pairing form statement \ref{pf:open0a:B3} holds for $v$.
\formallabel{GlobalTheory.T.pairForm_of_conjVec_of_botHyp}
\end{proposition}

\begin{proof}
The realisation carries an \emph{evaluation tie}: a vector $v_0$ of the models with $\|v_{0,\tau}\|=1$ for all $\tau$ such
that, for every $T$ in the good part at the true level, the translated evaluation at $t_S$ of $T(\otimes v_0)$ is almost
everywhere a function $F$ with $\bar F\in H_X$, and $s_{EE}$ maps the corresponding element of $E$ to $\bar F$
(\lref{GlobalTheory.MiTop.PairingL2.exists_eval}). We claim that there is $l\ne0$ with
$\mathrm{rf}(e)(\otimes v_0)=l\,\mathrm{rf}(e)(\otimes v)$ for all $e$ in the true part. If $\otimes v_0=l\otimes v$ for
some $l\ne0$ this is clear. Otherwise take $l=1$; the identity is trivial when $\mathrm{rf}(e)=0$, and if
$\mathrm{rf}(e)\ne0$ then $\mathrm{rf}(e)$ is a non-zero element of $\openzeroaRow$ (\ref{pf:open0a:B1}) whose
evaluation at $\otimes v_0$ is, by the evaluation tie, conjugate to a cusp form, and
\lref{GlobalTheory.T.unit_vector_rows_conjugate_alt} (by \ref{pf:open0a:P2} and \ref{pf:open0a:P12}: a unit vector
with this property is a non-zero multiple of the conjugate vector) gives $\otimes v_0=l'\otimes v$ with $l'\ne0$,
contrary to the assumption. Let $u_S$ and $\kappa_S>0$ be the normalising constants of the realisation ($|u_S|=1$) and put
$\kappa=\overline{u_S\kappa_S}\,l\ne0$. For $e$, $u$, $m$, $x$, $C$ and $b$ as in \ref{pf:open0a:B3}, the inner product
$\langle s_{EE}(e),P_m(\openzeroacj(x))\rangle$ is computed through the join of the realisation
(\lref{GlobalTheory.MiTop.PairingL2.inner_join}) as $\openzeroacj$ of an inner product in $H_X$ between the image of
$e$ and the period function of the ramified coordinates of $x$; the latter is $\varpi_C$
(\lref{GlobalTheory.MiTop.PairingL2.pwOf_op_ramCrd_eq_puFun}, using the bottom reading of the pins), and the former is
$u_S\kappa_S$ times the conjugate of the translated evaluation of $\mathrm{rf}(e)(\otimes v_0)=l\,\mathrm{rf}(e)(\otimes v)$.
By the definition of $b$ and \lref{GlobalTheory.MiTop.PairingL2.inner_smul_star}, the result is
$\openzeroacj(\kappa b)$.
\end{proof}

\begin{proposition}\label{pf:open0a:botrel}
Assume \ref{pf:open0a:P2}, \ref{pf:open0a:P12} and the bottom hypotheses for $W$. For $z\in G_f\cap Z(G)$ and a bottom
datum $d$ pinned at $z$, the bottom relation $\mathrm{BR}(z,d)$ holds (statement \ref{pf:open0a:B5}).
\formallabel{GlobalTheory.T.botRel_of_botHyp_of_pin_central}
\end{proposition}

\begin{proof}
Put $g_0=\iota(1,x_u)\,z\,t_S$ and $\mathrm{Bes}=\mathrm{Bes}_{d.B}$. We use:
\begin{enumerate}[label=(\alph*)]
\item (\lref{GlobalTheory.T.exists_isometry_colForm_eq_rTrans_of_pin}) a linear isometry $R$ of $M_u$ such that, if
$C$ is a column form of $R\varphi$ and $C_0$ a column form of $\varphi$, then the Bessel column forms satisfy
$\mathrm{cf}_{C_0}(\xi)=R_{d.a^{-1}}\,\mathrm{cf}_C(\xi)$ for all $\xi$;
\item (\lref{GlobalTheory.T.dictionaryColumnScalarForm_eq_traceColumnForm}) vectors $\xi(y)$, indexed by the reading
points $y$ and spanning the representation space of the Bessel datum, such that the reading at level $0$ of the column with
dictionary image $C$ is $g\mapsto\mathrm{cf}_C(\xi(y))(g\,d.a_0^{-1})$;
\item (\lref{GlobalTheory.T.conjugateRow_pairing_poincareSeries_besselPeriods}; about $1.1\cdot10^3$ lines) the
\emph{period law}: for every column form $C$, every $v$ and every $\phi\in L^2_X$ representing a translated evaluation of a
row, $\langle\bar\phi,\varpi_C\rangle$ is an integral over the archimedean group of the Hida setting of a fixed kernel $c(y)$
times Bessel periods of $\mathrm{cf}_C(W_u(y))$ against the corresponding right translates of the row;
\item the multiplicity maps $T\in\openzeroaRow$ satisfy $T(v)(gk)=\chi(k)\,T(v)(g)$ for $k\in J$, and $\chi(k)\ne0$ on $J$;
right translation by an archimedean element is realised on $\openzeroaRow$ by an operator $\sigma$ of the smooth tensors
(\lref{GlobalTheory.RealTransSm.rTrans_eq_map}); and the Bessel periods satisfy
$\mathrm{Bes}(f(\cdot\,h^{-1}),f')=\mathrm{Bes}(f,R_{\mathrm{emb}(h)}f')$ for $h\in\Uni(X_0)(\A)$.
\end{enumerate}
\emph{Step 1.} For $\varphi\in M_u$, a column form $C$ of $R\varphi$, $T\in\openzeroaRow$ and all $y,v$:
\begin{equation}\label{pf:open0a:eq:L1}
d.\mathrm{per}(\varphi,T)(y,v)=\chi(d.k)\,\mathrm{Bes}\bigl(\mathrm{cf}_C(\xi(y)),\,R_{g_0}(T(v))\bigr).
\end{equation}
Indeed, by (b) and (a) (with $C_0$ the column of $\varphi$ itself) and the centrality of $d.a$, the reading of
$e_u\varphi$ at $y$ is $g\mapsto\mathrm{cf}_C(\xi(y))(g\,(d.a_0d.a)^{-1})$. By (d), the period of this function against
$R_{d.b}(T(v))$ equals $\mathrm{Bes}(\mathrm{cf}_C(\xi(y)),R_{\mathrm{emb}(d.a_0d.a)d.b}(T(v)))$; by the pin,
$\mathrm{emb}(d.a_0d.a)\,d.b=g_0\,d.k$, and $T(v)(g\,g_0\,d.k)=\chi(d.k)\,T(v)(g\,g_0)$.

\emph{Step 2.} Let $\varphi,\varphi',T,T',\alpha,\alpha'$ be as in the definition of $\mathrm{BR}(z,d)$, with
$\alpha\,d.\mathrm{per}(\varphi,T)+\alpha'\,d.\mathrm{per}(\varphi',T')=0$, and let $C,C'$ be column forms of $R\varphi$,
$R\varphi'$. By \eqref{pf:open0a:eq:L1} and $\chi(d.k)\ne0$,
$\alpha\,\mathrm{Bes}(\mathrm{cf}_C(\xi(y)),R_{g_0}T(v))+\alpha'\,\mathrm{Bes}(\mathrm{cf}_{C'}(\xi(y)),R_{g_0}T'(v))=0$ for
all $y,v$. Since the $\xi(y)$ span and $\mathrm{cf}_C$ and the Bessel periods are linear, the same holds with $\xi(y)$
replaced by any vector $\xi'$.

\emph{Step 3.} Let $b,b'$ be pairings of $T(v),T'(v)$ at $zt_S$ with $C,C'$. Since the level matches at $zt_S$, the
translated evaluations are right translates by $g_0$ of rows, the archimedean translations being realised by $\sigma$
(d); by the period law (c), $\alpha b+\alpha'b'$ is the integral of $c(y)$ times
$\alpha\,\mathrm{Bes}(\mathrm{cf}_C(W_u(y)),R_{g_0}T(\sigma_yv))+\alpha'\,\mathrm{Bes}(\mathrm{cf}_{C'}(W_u(y)),R_{g_0}T'(\sigma_yv))$,
which vanishes for every $y$ by Step 2. Hence $\alpha b+\alpha'b'=0$.
\end{proof}

%==============================================================================
\subsection{Remarks on the formal statements}\label{sec:pf:open0a:remarks}

\begin{remark}[The hypotheses of the main theorem]\label{pf:open0a:rem:hyps}
(a) The hypothesis that the reading data of every weighted point admit a reading structure is used only in
\cref{pf:open0a:truepack}, through the leading reads. In the application in the proof of \cref{thm:families:open} the
parameter is of extension type, and the hypothesis is supplied by a component of the witness of the point.

(b) The diagonal row hypothesis is used only in \cref{pf:open0a:punz}. In the application it is supplied by the
convergent Siegel--Weil formula through \lref{GlobalTheory.T.diagonalRow_of_siegelWeil_unitary}, which derives
$(\mathrm{DR}_{E_{\mathrm{diag}}})$ from the Siegel--Weil formula of \cref{lit:families:sw}.

(c) Each of \ref{pf:open0a:P1}--\ref{pf:open0a:P13} is passed to at least one ingredient: \ref{pf:open0a:P1} and
\ref{pf:open0a:P5} to the point conditions and to the surjectivity along chains; \ref{pf:open0a:P2}, \ref{pf:open0a:P3}
and \ref{pf:open0a:P4} moreover to the family conditions, and \ref{pf:open0a:P2}, \ref{pf:open0a:P4} to the final
construction; \ref{pf:open0a:P6}--\ref{pf:open0a:P11} only to the surjectivity along chains (\cref{pf:open0a:onto});
\ref{pf:open0a:P12} only to the final construction (\cref{pf:open0a:mistarpunz}); \ref{pf:open0a:P13} only to the family
conditions (\cref{pf:open0a:famall}). Whether each input is actually used inside the cited ingredients has not been
checked here.
\end{remark}

\begin{remark}[Nature of the conclusion]\label{pf:open0a:rem:vacuity}
The main theorem is an eventual statement about bundles \emph{linked} to good points: it asserts nothing about the
existence of such bundles, which is the subject of \cref{thm:member-bundles}. Linkage requires a bundle witness, which
already contains much of the data from which the point conditions and the truth of the real package are read off
(\cref{sec:pf:open0a:points}). \Cref{pf:open0a:mitrue} is a pointwise statement: it holds for every true real package
with $r_S\ge1$ and non-vanishing periods, without any size assumption; the only eventual inputs are the size
conditions used in \cref{pf:open0a:points,pf:open0a:punz}.
\end{remark}

\begin{remark}[Unused hypotheses and the role of the constants]\label{pf:open0a:rem:unused}
(a) \Cref{pf:open0a:served} holds for every real package; its formal statement carries the bottom hypotheses and the
existence of a recursion as hypotheses, which are not used.
(b) In \cref{pf:open0a:mitrue} the hypothesis $r_S\ge1$ is used only through \ref{pf:open0a:F1}.
(c) The function arguments $F,G,H$ of \ref{pf:open0a:S1} are instantiated by constants ($0$, except
$G\equiv\max(1801,C^{\mathrm{nc}}_{l,0})$ in the proof of \cref{pf:open0a:points}); the dependence on $(C_g,C_T)$ is
exploited only through \ref{pf:open0a:S2}.
(d) Of the standing size conditions, the proof of \cref{pf:open0a:punzsdiag} uses only the inequalities
\eqref{pf:open0a:eq:HDused}; the other components are needed to apply the diagonal row hypothesis.
\end{remark}

\begin{remark}[The non-vanishing argument]\label{pf:open0a:rem:punz}
The non-vanishing of periods (\cref{pf:open0a:punzsdiag}) is a dimension count: if all ramified vectors of a word had
vanishing period, the kernel of $P_u$ would have dimension at least $d_U(u)\ge(1-f/(T_u-1))\,d_0(u)\ge\frac34d_0(u)$,
whereas the diagonal bound gives $\dim\ker P_u\le\openzeroahs(u)\le\openzeroaeD\,d_0(u)$ with $\openzeroaeD<\frac18$.
It is used in the final construction only once per word, to produce one non-zero pairing
(\cref{pf:open0a:anchorpairs}(c)); all other pairings are then controlled by linear algebra
(\cref{pf:open0a:anchorword}).
\end{remark}
